% Begin preamble
\documentclass[11pt]{article}
\usepackage[T1]{fontenc}
\usepackage[utf8]{inputenc}
\usepackage{lmodern}
\usepackage[letterpaper,margin=0.95in]{geometry}
\usepackage{microtype,amsmath,amssymb,amsthm,mathtools,bm}
\usepackage{booktabs,longtable,array,enumitem,xcolor}
\usepackage{fancyhdr}
\usepackage{hyperref}
\definecolor{linkblue}{RGB}{26,65,108}
\hypersetup{colorlinks=true,linkcolor=linkblue,citecolor=linkblue,urlcolor=linkblue,
 pdftitle={Higher Reflection Identities and Resonant Zeta Calculus},
 pdfauthor={ProveIt research continuation},pdfsubject={Exact identities and rigorous research continuation}}
\pagestyle{fancy}
\fancyhf{}
\fancyhead[L]{\small Higher Reflection Identities}
\fancyhead[R]{\small ProveIt research continuation}
\fancyfoot[C]{\thepage}
\setlength{\headheight}{14pt}
\setlength{\emergencystretch}{2em}
\setlength{\parskip}{3pt}
\setlist{itemsep=3pt,topsep=5pt}
\numberwithin{equation}{section}
\newtheorem{theorem}{Theorem}[section]
\newtheorem{proposition}[theorem]{Proposition}
\newtheorem{lemma}[theorem]{Lemma}
\newtheorem{corollary}[theorem]{Corollary}
\theoremstyle{definition}
\newtheorem{definition}[theorem]{Definition}
\newtheorem{question}[theorem]{Research question}
\newtheorem{example}[theorem]{Example}
\theoremstyle{remark}
\newtheorem{remark}[theorem]{Remark}
\DeclareMathOperator{\FP}{FP}
\DeclareMathOperator{\CT}{CT}
\DeclareMathOperator{\Li}{Li}
\DeclareMathOperator{\Res}{Res}
\DeclareMathOperator{\Rea}{Re}
\newcommand{\dd}{\,\mathrm d}
\newcommand{\C}{\mathbb C}
\newcommand{\Q}{\mathbb Q}
\newcommand{\Z}{\mathbb Z}
\newcommand{\N}{\mathbb N}
\newcommand{\src}[1]{\texttt{\detokenize{#1}}}
\newcommand{\ind}{\mathbf 1}


\begin{document}
\title{\textbf{Higher Reflection Identities\\and Resonant Zeta Calculus}\\[0.6em]
\large Complex-power sine kernels, harmonic parity,\\
geometric collisions, and quartic Tornheim derivatives}
\author{Research continuation for the ProveIt project}
\date{11 October 2026\quad (UTC)}
\maketitle
\begin{abstract}
We develop four exact continuations of the current ProveIt polylogarithm
manuscript and its newest incoming reports. A complex-power sine weight
produces a two-parameter Hurwitz transform with explicit Gamma Fourier
coefficients and a polylogarithmic Euler integral. Removing its complete
crossing polar term gives every finite-part moment of an argument
derivative of a generalized Stieltjes function against a logarithmic
sine power. The first nonlinear specialization identifies a combination
of the digamma cube and a logarithmic sine-square moment in ordinary
zeta and Stieltjes constants. A second result characterizes exactly the
polynomials in generalized harmonic numbers whose centered Dirichlet
series have no poles at any nonpositive even integer. A third gives an
explicit analytic subtraction for arbitrary finite clusters of periodic
digamma factors, including hierarchical collisions. Finally, we compute
the complete fourth directional derivative of Tornheim zeta, give
convergent representatives for its three formal coordinates, and prove
an all-order dimension theorem for the specified symmetry and boundary
relations. It explains why one diagonal coordinate suffices for cyclic
cubic and quartic sums and why the same argument fails at order five.
All statements include their continuation or finite-part conventions.
The package contains reproducible exact algebra and independent numerical
diagnostics. The Gaussian $S_6$ and revised $S_8$ evaluations remain open.
\end{abstract}
\clearpage
\begingroup
\small
\setlength{\parskip}{0pt}
\tableofcontents
\endgroup
\clearpage
% Begin sections/01_scope
\section{Scope, provenance, and conventions}
\label{sec:scope}

The starting point is the ProveIt repository at revision
\begin{center}
\texttt{c2cb285b64273e4cb2c8d7a03c68a4a99ac812d4}.
\end{center}
The source inventory under
\src{Analysis/Polylogarithms/docs/manuscript} and all nine ZIP reports
present under \src{docs/incoming} at that revision were retrieved.
Selected canonical manuscript sections were read alongside the reports.
The focused reading concerned the integration, harmonic, reflection,
Stieltjes contact, and Tornheim sections. This is a continuation of those
parts, not a claim to have audited every proof in the four-volume project.
The accompanying \src{SOURCE_AUDIT.md} identifies the relevant files and
the exact boundary between inherited results and the present extensions.

The latest incoming reports matter mathematically. They already prove
the cubic digamma--harmonic bridge, the independent cubic Tornheim
formula, the polynomial reflected Stieltjes closure, and several
all-depth spectral completion results. Repeating any of these as a new
theorem would misrepresent the state of the project. The questions
continued here are more specific:
\begin{enumerate}
\item The nonpolynomial reflected generator requested in Q8 of
\emph{Exact Resonant Identities} \cite{RepoExact}.
\item The necessary and sufficient centered cancellation criterion for
generalized harmonic polynomials requested in Q11 of the same report.
\item Simultaneous and hierarchical geometric collisions beyond the
two-factor case, asked for in the periodic contact reports
\cite{RepoPeriodic,RepoMixed}.
\item The quartic directional layer and its cancellation map, requested
in \emph{Independent Orders and Exact Reductions}, Section 7,
and the cubic harmonic report \cite{RepoIndependent,RepoCubic}.
\end{enumerate}

These were posed as research questions. We do not recast them as named
classical conjectures. The results below resolve the stated analytic and
formal-algebraic components; the arithmetic reduction of the retained
constants is a separate problem.

\subsection{Basic functions and the meaning of a derivative}

We use
\[
 \zeta(s,a)=\sum_{n=0}^{\infty}(n+a)^{-s},\qquad
 \Li_s(z)=\sum_{n=1}^{\infty}\frac{z^n}{n^s}
\]
in their initial convergence domains, followed by the explicitly stated
analytic continuation. The generalized Stieltjes functions are defined by
\begin{equation}\label{scope:stieltjes}
 \zeta(1+u,x)=\frac1u+
       \sum_{m=0}^{\infty}\frac{(-1)^m}{m!}\gamma_m(x)u^m.
\end{equation}
In particular $\gamma_0(x)=-\psi(x)$, where
$\psi=\Gamma'/\Gamma$. Write $\gamma=\gamma_0(1)$ and
$\gamma_m=\gamma_m(1)$. A prime on $\psi$ or $\gamma_m$ means an
argument derivative. A superscript $\zeta^{(j)}(s)$ means an order
derivative of the Riemann zeta function. Multivariable subscripts on
the Tornheim remainder mean partial derivatives, always evaluated at
the point specified. These distinctions are essential at a spectral pole.

The generalized harmonic numbers are
\[
 H_n^{(r)}=\sum_{k=1}^n k^{-r},\qquad H_n=H_n^{(1)}.
\]
All powers of positive real variables use the real logarithm. Thus, for
complex $\lambda$, the expression $(\sin(\pi x)/\pi)^{2\lambda}$
on $0<x<1$ is unambiguous. Meromorphic identities are continued from
a nonempty open convergence domain; a value at a crossing of polar
divisors is never assigned by informal substitution.

\subsection{Coordinate finite parts}

For an endpoint expansion in real powers and integral logarithmic powers,
the unit-coordinate finite part is the coefficient independent of both
$\epsilon$ and $\log\epsilon$ in the cutoff integral. On the unit interval,
\begin{equation}\label{scope:FP}
 \FP\int_0^1 f(x)\dd x
   =\CT_{\epsilon\downarrow0}\int_\epsilon^{1-\epsilon}f(x)\dd x.
\end{equation}
The coordinates are $x$ at zero and $1-x$ at one. Independent endpoint
cutoffs give the same additive constant. For example,
\[
 \FP\int_0^1x^{-1}\log^j x\dd x=0,\qquad
 \FP\int_0^1x^{-r}\dd x=\frac1{1-r}\quad(r\ne1).
\]
At periodic seams, the same convention is applied on each side in the
unit argument coordinate. Geometric collision constants are then taken
in a second parameter after this first finite part has been formed.
The two operations can differ from first merging the factors.

All coefficient extractions below are from a germ proved holomorphic
after its full prescribed subtraction. Numerical quadrature is used
only to test identities already supplied with proofs. Exact finite
linear algebra establishes a statement about its specified relation
space, not arithmetic independence of the associated periods.

\subsection{Relation to the literature}

Hurwitz Fourier transforms, Beta and Gamma identities, and their
applications to logarithmic integrals are classical
\cite{DLMFHurwitz,DLMFBeta,EspinosaMoll}. Harmonic Stieltjes constants
and log-sine connections were studied before this project
\cite{Kargin,Can}. The Mellin continuation and Euler boundary relations
for Tornheim zeta are likewise established tools \cite{BorweinDilcher}.
The contributions claimed here are the particular completed generators,
the exact all-polynomial criterion, the explicit cluster subtraction,
and the higher directional reconstruction proved below. The source
search was focused; it does not establish worldwide priority for every
displayed specialization.


% Begin sections/02_beta
\section{A completed complex-power sine transform}
\label{sec:beta}

Put
\begin{equation}\label{beta:weight}
 \ell(x)=\log\frac{\sin\pi x}{\pi},\qquad L=\log(2\pi),\qquad
 W_\lambda(x)=e^{2\lambda\ell(x)}.
\end{equation}
The normalization by $\pi$ makes the leading endpoint coefficient equal
to one. It is useful both for coordinate finite parts and for keeping
the later crossing subtraction unambiguous. Define
\begin{equation}\label{beta:A}
 A(s,\lambda)=\int_0^1\zeta(s,x)W_\lambda(x)\dd x,
 \qquad
 C_0(\lambda)=\frac{\Gamma(1+2\lambda)}
 {\Gamma(1+\lambda)^2(2\pi)^{2\lambda}}.
\end{equation}
An initial domain sufficient for everything in the first theorem is
$\Re\lambda>-1/2$, $\Re s<0$. All subsequent values of $A$ mean its
meromorphic continuation.

\subsection{Fourier coefficients and a polylogarithmic Euler integral}

\begin{theorem}[The complex-power transform]\label{beta:transform}
For $\Re\lambda>-1/2$ the Fourier coefficients of $W_\lambda$ are
\begin{equation}\label{beta:cn}
 c_n(\lambda)=\int_0^1W_\lambda(x)e^{-2\pi inx}\dd x
 =\frac{(-1)^n\Gamma(1+2\lambda)}
 {(2\pi)^{2\lambda}\Gamma(1+\lambda+n)\Gamma(1+\lambda-n)}.
\end{equation}
They satisfy $c_{-n}=c_n$, $c_0=C_0$, and, for $n\geq1$,
\begin{equation}\label{beta:pochhammer}
 c_n(\lambda)=C_0(\lambda)\frac{(-\lambda)_n}{(1+\lambda)_n}.
\end{equation}
Consequently,
\begin{equation}\label{beta:dirichlet}
 A(s,\lambda)=B(s)\sum_{n=1}^{\infty}c_n(\lambda)n^{s-1},\qquad
 B(s)=2\Gamma(1-s)(2\pi)^{s-1}\sin\frac{\pi s}{2}.
\end{equation}
In the strip $-1/2<\Re\lambda<0$, $\Re s<0$, a second representation is
\begin{align}\label{beta:euler}
 A(s,\lambda)
 &=-\frac{2}{\pi}\Gamma(1-s)(2\pi)^{s-1-2\lambda}
      \sin\frac{\pi s}{2}\sin(\pi\lambda)\notag\\
 &\hspace{2em}\times
 \int_0^1t^{-\lambda-1}(1-t)^{2\lambda}\Li_{1-s}(t)\dd t.
\end{align}
Both formulas continue meromorphically wherever their two sides are
defined by continuation.
\end{theorem}

\begin{proof}
Reflection about $1/2$ kills the sine coefficient. If
$I_n=\int_0^\pi\sin^{2\lambda}t\cos(2nt)\dd t$, integration of
$\frac{\mathrm d}{\mathrm dt}
 [\sin^{2\lambda+1}t\cos((2n+1)t)]$ gives
\[
 (\lambda-n)I_n+(\lambda+n+1)I_{n+1}=0.
\]
The boundary term vanishes for $\Re\lambda>-1/2$.
The beta integral gives
$I_0=\sqrt\pi\,\Gamma(\lambda+1/2)/\Gamma(\lambda+1)$.
The recurrence and duplication formula prove
\eqref{beta:cn}--\eqref{beta:pochhammer}. The apparent exceptions at
integer parameters are read by continuity.

The Hurwitz Fourier expansion for $\Re s<0$ is absolutely convergent.
Its sine terms integrate to zero and its cosine terms give
\eqref{beta:dirichlet}; see also \cite[\S25.11]{DLMFHurwitz}.
The decay $c_n=O(n^{-1-2\Re\lambda})$ shows that the resulting
Dirichlet series itself converges absolutely for
$\Re s<1+2\Re\lambda$, away from exceptional parameters interpreted
by continuity. Lastly, Gamma reflection gives
\[
 c_n(\lambda)=-\frac{\sin\pi\lambda}{\pi(2\pi)^{2\lambda}}
 \int_0^1t^{n-\lambda-1}(1-t)^{2\lambda}\dd t.
\]
Absolute convergence in the stated strip permits summation under the
integral, producing \eqref{beta:euler}.
\end{proof}

These Fourier and beta mechanisms are classical. In particular,
Espinosa and Moll treat Hurwitz transforms of trigonometric powers
and the first log-sine transform \cite[\S\S5,13]{EspinosaMoll}.
The purpose of Theorem~\ref{beta:transform} here is to supply a common
generator whose \emph{completed germ at the joint resonance} can be
specified at every Stieltjes and argument-derivative order.

For an integer $m\geq1$ the Fourier expansion terminates:
\begin{equation}\label{beta:integer}
 A(s,m)=\frac{B(s)}{(2\pi)^{2m}}
 \sum_{n=1}^m(-1)^n\binom{2m}{m+n}n^{s-1}.
\end{equation}
For example, comparison of the constant term at $s=1$ gives the
ordinary convergent identity
\begin{equation}\label{beta:psiinteger}
 \int_0^1\psi(x)W_m(x)\dd x
 =-C_0(m)(\gamma+L)
 +\frac{2}{(2\pi)^{2m}}
   \sum_{n=1}^m(-1)^n\binom{2m}{m+n}\log n.
\end{equation}
Higher Laurent coefficients give the analogous all-index Stieltjes
moments as finite Gamma and logarithmic expressions.

\subsection{The complete crossing polar term}

Let $P_p(u)=(1+u)_p$, with $P_0=1$, and define the entire-in-$\lambda$
Taylor coefficients
\begin{align}\label{beta:d}
 d_j(\lambda)
 &=[x^{2j}]\left(\frac{\sin\pi x}{\pi x}\right)^{2\lambda}\notag\\
 &=[z^j]\exp\left(-2\lambda
                 \sum_{r=1}^{\infty}\frac{\zeta(2r)}r z^r\right).
\end{align}
Thus $d_0=1$, $d_1=-2\lambda\zeta(2)$, and
$d_2=-\lambda\zeta(4)+2\lambda^2\zeta(2)^2$.

\begin{theorem}[All logarithmic Stieltjes moments]\label{beta:completion}
For every integer $p\geq0$ the expression
\begin{align}\label{beta:Hp}
 \mathcal H_p(u,\lambda)
 ={}&(-1)^pP_p(u)A(1+p+u,\lambda)
       -\ind_{p=0}\frac{C_0(\lambda)}u\notag\\
 &-\ind_{2\mid p}
       \frac{(-1)^pP_p(u)d_{p/2}(\lambda)}{2\lambda-u}
\end{align}
extends holomorphically to a neighborhood of $(u,\lambda)=(0,0)$.
For all integers $m,k\geq0$,
\begin{equation}\label{beta:allmoments}
 \boxed{\displaystyle
 \FP\int_0^1\gamma_m^{(p)}(x)\ell(x)^k\dd x
  =\frac{(-1)^m m!k!}{2^k}
      [u^m\lambda^k]\mathcal H_p(u,\lambda).}
\end{equation}
The finite part is exactly the unit-coordinate convention
\eqref{scope:FP}.
\end{theorem}

\begin{proof}
The Hurwitz shift equation gives, near $x=0$,
\begin{equation}\label{beta:local}
 \partial_x^p\left(\zeta(1+u,x)-\frac1u\right)
 =(-1)^pP_p(u)x^{-1-p-u}+h_p(u,x),
\end{equation}
where $h_p$ is holomorphic jointly near $(u,x)=(0,0)$.
For $p=0$ the subtraction $1/u$ is retained inside $h_0$;
for $p>0$ its derivative is zero. Moreover
\[
 W_\lambda(x)=x^{2\lambda}\sum_{j\geq0}d_j(\lambda)x^{2j}.
\]
Integration of the singular product from zero to a fixed
$\delta\in(0,1)$ produces the denominators
$2\lambda-u+2j-p$. Only $j=p/2$ can vanish at the origin, and it
occurs only when $p$ is even. Its full numerator is the one displayed
in \eqref{beta:Hp}. Subtracting finitely many endpoint terms makes
the remaining integral locally normally convergent; all the other
denominators are bounded away from zero. The endpoint $x=1$ has only
the integrable analytic family $(1-x)^{2\lambda}$ times analytic
functions and creates no divisor through $(0,0)$. The pole $1/u$ of
the original Hurwitz function contributes $C_0(\lambda)/u$ only
when $p=0$. This proves joint holomorphy.

For completeness, the local resonant integral is
$N(u,\lambda)\delta^{2\lambda-u}/(2\lambda-u)$. After subtracting
the entire quotient $N(u,\lambda)/(2\lambda-u)$, its Taylor
coefficients are those of
$N(u,\lambda)(\delta^{2\lambda-u}-1)/(2\lambda-u)$.
These are precisely the constants left by integrating the corresponding
powers of $\log x$ in the coordinate $x$. All nonresonant terms have
the same property by ordinary analytic continuation of
$(\delta^a-\epsilon^a)/a$. Thus Taylor coefficients commute with the
specified finite part after the full subtraction. Applying
\eqref{scope:stieltjes} and
$W_\lambda=\sum_k2^k\ell^k\lambda^k/k!$ proves
\eqref{beta:allmoments}.
\end{proof}

\begin{corollary}\label{beta:zero}
For every $m,p\geq0$,
\[
 \FP\int_0^1\gamma_m^{(p)}(x)\dd x=0.
\]
\end{corollary}
\begin{proof}
The meromorphic identity $A(s,0)=0$ follows from
\eqref{beta:dirichlet}, since all nonzero Fourier coefficients
vanish. Also $d_j(0)=0$ for $j>0$. In \eqref{beta:Hp} the two
remaining terms for $p=0$ cancel along $\lambda=0$; for $p>0$ each
term vanishes there. Hence $\mathcal H_p(u,0)=0$ as a holomorphic
germ.
\end{proof}

\begin{remark}[A removable limit need not equal a coordinate finite part]
\label{beta:anomaly}
For an even integer $p\geq2$, meromorphic continuation from a
convergence domain gives
\begin{equation}\label{beta:anomalyformula}
 \lim_{\lambda\to0}
 \left(\int_0^1\psi^{(p)}(x)W_\lambda(x)\dd x\right)_{\mathrm{cont}}
 =2(p-1)!\zeta(p),
 \qquad
 \FP\int_0^1\psi^{(p)}(x)\dd x=0.
\end{equation}
Indeed $d_{p/2}'(0)=-4\zeta(p)/p$, and the unsubtracted crossing
quotient along $u=0$ is
$-p!d_{p/2}(\lambda)/(2\lambda)$ for the polygamma integral.
This gives the first value, while Corollary~\ref{beta:zero} gives the
second. For $p=2$ the discrepancy is $2\zeta(2)$. The correction is
to remove the \emph{full analytic numerator}, including terms that
vanish at the crossing. This is a normalization warning about a
tempting extension, not a claim that the incoming completed-germ
theorems make this error.
\end{remark}

\subsection{Cotangent powers and normalized antiderivatives}

Let $K(x)=\pi\cot\pi x$. The identity
$K^2=W_{\lambda-1}/W_\lambda-\pi^2$ implies
\begin{equation}\label{beta:evencot}
 \int_0^1\zeta(s,x)K(x)^{2h}W_\lambda(x)\dd x
 =\sum_{j=0}^h\binom hj(-\pi^2)^{h-j}A(s,\lambda-j).
\end{equation}
Since $W_\mu'=2\mu K W_\mu$ and
$\partial_x\zeta(s,x)=-s\zeta(s+1,x)$, integration by parts in
an endpoint-convergent domain gives
\begin{equation}\label{beta:oddcot}
 \int_0^1\zeta(s,x)K(x)^{2h+1}W_\lambda(x)\dd x
 =\sum_{j=0}^h\binom hj(-\pi^2)^{h-j}
   \frac{s}{2(\lambda-j)}A(s+1,\lambda-j).
\end{equation}
These are meromorphic identities. A zero denominator on the right
requires continuation and, for a coordinate finite part, the same
local completion as above. The polynomial cotangent closure already
proved in \cite{RepoExact} is therefore embedded in a complex-power
family rather than claimed afresh.

The symmetric Gamma primitive is particularly simple:
\begin{proposition}\label{beta:gammaprimitive}
For $\Re\lambda>-1/2$,
\begin{equation}\label{beta:gamma}
 \int_0^1\log\Gamma(x)W_\lambda(x)\dd x=-\frac14C_0'(\lambda).
\end{equation}
Consequently every $k\geq0$ satisfies the ordinarily convergent identity
\begin{equation}\label{beta:gammalogs}
 \int_0^1\log\Gamma(x)\ell(x)^k\dd x
 =-\frac{C_0^{(k+1)}(0)}{2^{k+2}}.
\end{equation}
\end{proposition}
\begin{proof}
Gamma reflection gives
$\log\Gamma(x)+\log\Gamma(1-x)=-\ell(x)$. Multiply by the symmetric
weight and integrate. The right side is $-C_0'(\lambda)/2$, while
the left side is twice the desired integral. Differentiation is
justified by endpoint domination on compact subsets of the stated
half-plane.
\end{proof}

For arbitrary Stieltjes index, the normalized antiderivative generator
is
\begin{equation}\label{beta:primitivegenerator}
 \mathcal P(u,x)=-\frac{\zeta(u,x)-\zeta(u,1)+x-1}{u}.
\end{equation}
Its apparent singularity at $u=0$ is removable because
$\zeta(0,x)=1/2-x$. The Hurwitz derivative identity proves
\begin{equation}\label{beta:primitivecoeff}
 \partial_x\mathcal P(u,x)=\zeta(1+u,x)-\frac1u,
 \qquad
 \int_1^x\gamma_m(t)\dd t=(-1)^m m![u^m]\mathcal P(u,x).
\end{equation}
In particular $\mathcal P(0,x)=-\log\Gamma(x)$ by Lerch's formula.
The normalization at $x=1$ fixes the integration constant before any
reflection or finite-part operation. Iterating ordinary integration
of this holomorphic generator gives normalized higher primitives.

\subsection{Harmonic jets of the weight}

\begin{proposition}\label{beta:harmonicjets}
For every integer $n\geq1$ the germ of its Fourier coefficient is
\begin{align}\label{beta:jet}
 c_n(\lambda)=-\frac\lambda n\exp\bigg(&-2L\lambda
 +\sum_{j\geq2}\frac{(-1)^j(2^j-2)\zeta(j)}j\lambda^j\notag\\
 &-\sum_{j\geq1}
 \frac{H_{n-1}^{(j)}+(-1)^{j+1}H_n^{(j)}}j\lambda^j\bigg).
\end{align}
Thus every logarithmic sine moment is generated by finite polynomials
in generalized harmonic numbers and ordinary zeta constants.
\end{proposition}
\begin{proof}
Factor \eqref{beta:pochhammer} as
\[
 c_n=-\frac\lambda n C_0(\lambda)
  \frac{\prod_{j=1}^{n-1}(1-\lambda/j)}
       {\prod_{j=1}^{n}(1+\lambda/j)}.
\]
Expanding the logarithm of each finite product and using the Taylor
series of $\log\Gamma(1+\lambda)$ gives \eqref{beta:jet}.
Each fixed coefficient requires only finitely many terms. Initially
the differentiated Dirichlet series converges normally in a smaller
half-plane; its extension is then supplied by
Theorem~\ref{beta:completion}.
\end{proof}

The first two terms illustrate the inclusive harmonic convention:
\begin{equation}\label{beta:firstjets}
 c_n(\lambda)=-\frac\lambda n+\lambda^2
 \left(\frac{2(L+H_n)}n-\frac1{n^2}\right)+O(\lambda^3).
\end{equation}
More explicitly, if $J_k(s)=\int_0^1\zeta(s,x)\ell(x)^k\dd x$,
then in its initial domain
\begin{equation}\label{beta:belltransform}
 J_k(s)=\frac{k!}{2^k}B(s)
 \sum_{n\geq1}[\lambda^k]c_n(\lambda)n^{s-1}.
\end{equation}
This is a finite Bell-polynomial construction at each fixed $k$, not
an assertion that all resulting harmonic constants have already been
reduced to ordinary ones.

\subsection{The first two logarithmic moments and an exact cancellation}

For $k=1$, the classical transform becomes
\begin{align}\label{beta:J1}
 J_1(s)&=-\Gamma(1-s)(2\pi)^{s-1}
                \sin\frac{\pi s}{2}\zeta(2-s)\notag\\
 &=\frac{\pi\cot(\pi(s-1)/2)}{s-1}\zeta(s-1).
\end{align}
The second expression follows from the Riemann functional equation.
The completed germ is
\begin{equation}\label{beta:J1completion}
 J_1(1+u)+\frac1{u^2}+\frac L u.
\end{equation}
Its $u^m$ coefficient, multiplied by $(-1)^m m!$, is
$\FP\int\gamma_m\ell$. In particular,
\begin{equation}\label{beta:firstmoment}
 \FP\int_0^1\gamma_0(x)\ell(x)\dd x
 =\zeta''(0)+\frac{\zeta(2)}2.
\end{equation}
This also cross-checks the log-Gamma--cotangent specialization in
\cite{RepoExact}; neither \eqref{beta:J1} nor that specialization is
claimed to be new.

Write the \emph{inclusive} harmonic zeta function as
\begin{equation}\label{beta:Hzeta}
 \zeta_H(v)=\sum_{n\geq1}\frac{H_n}{n^v}
 =\frac1{(v-1)^2}+\frac\gamma{v-1}
  +\sum_{j\geq0}\frac{(-1)^j\gamma_H(j)}{j!}(v-1)^j.
\end{equation}
Here $\gamma_H(0)=(\gamma^2+\zeta(2))/2$; this normalization and
its Laurent theory are established in \cite{Kargin,Can}.

\begin{theorem}[The sine-square bridge]\label{beta:secondmoment}
One has
\begin{equation}\label{beta:J2}
 J_2(s)=B(s)\left\{L\zeta(2-s)+\zeta_H(2-s)
                       -\frac12\zeta(3-s)\right\}.
\end{equation}
Its completion at $s=1+u$ is
\begin{equation}\label{beta:J2completion}
 J_2(1+u)+\frac2{u^3}-\frac{L^2+\zeta(2)/2}{u}.
\end{equation}
In particular,
\begin{align}\label{beta:P2}
 \FP\int_0^1\psi(x)\ell(x)^2\dd x
 ={}&2\gamma_H(1)+\zeta'(2)+2L\gamma_1
    +\frac{\gamma^3}3+\gamma^2L\notag\\
   &-\frac{2L^3}3+\frac{2\zeta(3)}3.
\end{align}
Combining this with the already proved cubic harmonic bridge yields
the ordinary-constant identity
\begin{align}\label{beta:cancellation}
 \boxed{\begin{aligned}
 \FP\int_0^1\bigl\{\psi(x)^3-3\psi(x)\ell(x)^2\bigr\}\dd x
 ={}&3\zeta(2)+3\zeta'(2)-6(\gamma+L)\gamma_1\\
   &-\gamma^3-3\gamma^2L+2L^3-2\zeta(3).
 \end{aligned}}
\end{align}
\end{theorem}

\begin{proof}
Equation \eqref{beta:J2} is the second coefficient of
\eqref{beta:firstjets} in \eqref{beta:belltransform}.
The polar terms of \eqref{beta:J2completion} also follow directly
by taking the second $\lambda$ coefficient of
$C_0(\lambda)/u+1/(2\lambda-u)$ in
Theorem~\ref{beta:completion}. To make the constant calculation
explicit, set $a=\gamma+L$. Then
\[
 B(1+u)=-\frac2u
 \exp\left(au-\frac{\zeta(2)}4u^2+\frac{\zeta(3)}3u^3+O(u^4)\right),
\]
and the braces in \eqref{beta:J2} equal
\[
 \frac1{u^2}-\frac a u+L\gamma+\frac{\gamma^2}2
  +u\left(L\gamma_1+\gamma_H(1)+\frac{\zeta'(2)}2\right)+O(u^2).
\]
Multiplication gives the negative of \eqref{beta:P2}, because
$\gamma_0=-\psi$.

The incoming cubic report \cite{RepoCubic} proves, in precisely this
coordinate convention,
\begin{equation}\label{beta:inheritedcube}
 \FP\int_0^1\psi(x)^3\dd x
 =3\zeta(2)+6\gamma_H(1)+6\zeta'(2)-6\gamma\gamma_1.
\end{equation}
Subtracting three times \eqref{beta:P2} cancels $\gamma_H(1)$
identically and gives \eqref{beta:cancellation}. This use of the cubic
bridge is explicit: we prove the new cancellation, not a new proof
claim for the inherited bridge.
\end{proof}

The individual quantities are, for orientation,
\begin{align*}
 \FP\int_0^1\psi(x)^3\dd x
   &=1.9662627343915343406643753384116314\ldots,\\
 \FP\int_0^1\psi(x)\ell(x)^2\dd x
   &=-3.0645068706816349567044046185265270\ldots.
\end{align*}
Their combination in \eqref{beta:cancellation} is
$11.1597833464364392107775891939912123\ldots$.
These decimals come from independent endpoint-subtracted quadratures
and harmonic-tail evaluation in the supplied check script. They are
diagnostics; the coefficient identities and the inherited theorem are
the proof. No reduction of $\gamma_H(1)$ itself to ordinary constants
is inferred from its disappearance in this particular combination.


% Begin sections/03_harmonic
\section{A complete criterion for centered harmonic pole cancellation}
\label{hc:section}

Question~11 of the incoming Exact Resonant Identities report \cite{RepoExact} asks which
polynomials in generalized harmonic numbers have no principal part at
any nonpositive even spectral integer when the outer shift is $1/2$.
The criterion below is finite and necessary as well as sufficient.
The necessity is the substantive point: parity of the individual
asymptotic expansions alone supplies only one implication.

For a fixed $R\ge1$, write
\[
 H_n^{(r)}=\sum_{k=1}^n k^{-r},\qquad H_n^{(1)}=H_n,
 \qquad
 D_P(s)=\sum_{n=0}^{\infty}
 \frac{P(H_n^{(1)},\ldots,H_n^{(R)})}{(n+1/2)^s},
 \quad P\in\mathbb C[X_1,\ldots,X_R].
\]
The defining series is absolutely convergent for $\Re s>1$.
As usual $H_0^{(r)}=0$, so the finite $n=0$ term is retained.
Define an involution $\iota$ on the polynomial ring by
\begin{equation}\label{hc:involution}
 \iota(X_{2k})=2\zeta(2k)-X_{2k},\qquad
 \iota(X_{2k+1})=X_{2k+1},\qquad \iota(X_1)=X_1.
\end{equation}
Only variables whose indices are at most $R$ are used.

\begin{theorem}[Necessary and sufficient centered cancellation]
\label{hc:main}
The function $D_P$ has a meromorphic continuation to $\mathbb C$.
It is holomorphic at every point $s=0,-2,-4,\ldots$ if and only if
\begin{equation}\label{hc:criterion}
                         \iota(P)=P.
\end{equation}
Equivalently, put $Y_{2k}=X_{2k}-\zeta(2k)$. After expressing $P$ in
the variables $X_1,X_3,X_5,\ldots$ and $Y_2,Y_4,Y_6,\ldots$, every
monomial that occurs must have even total degree in the $Y$ variables.
The invariant algebra is therefore
\begin{equation}\label{hc:invariant-ring}
 \mathbb C\bigl[X_1,X_3,X_5,\ldots,
        Y_{2i}Y_{2j}\ (1\le i\le j\le\lfloor R/2\rfloor)\bigr].
\end{equation}
This is a generating description, not a claim that the displayed
quadratic generators are algebraically independent.
\end{theorem}

For example, every polynomial in odd-order harmonic numbers has the
cancellation, as does every product of an even number of even-order
tails. Thus the numerators
\[
 H_n^{(3)},\qquad
 H_n^2H_n^{(5)},\qquad
 (H_n^{(2)}-\zeta(2))^2,\qquad
 (H_n^{(2)}-\zeta(2))(H_n^{(4)}-\zeta(4))
\]
all satisfy the criterion. The two even-order factors need not be equal.
The raw numerator $(H_n^{(2)})^2$ fails: its simple-pole coefficient
at $s=0$ is $-2\zeta(2)$.

\subsection{Bernoulli coefficients and the principal parts}

Use $x=n+1/2$, $z=x^{-1}$, and the formal symbol $\mathcal L=\log x$. Define
\begin{align}
 h_1(z,\mathcal L)&=\mathcal L+\gamma-
           \sum_{k\ge1}\frac{B_{2k}(1/2)}{2k}z^{2k},\label{hc:h1}\\
 h_r(z)&=\zeta(r)-\frac{z^{r-1}}{r-1}
  -\sum_{k\ge1}\frac{B_{2k}(1/2)}{(2k)!}(r)_{2k-1}
                         z^{r+2k-1},\qquad r\ge2.\label{hc:hr}
\end{align}
These are formal asymptotic series; no convergence of the infinite
Bernoulli expansions is asserted. Each finite truncation is the
usual shifted polygamma or Hurwitz-zeta asymptotic expansion \cite{DLMFGamma,DLMFHurwitz}, with
a remainder sufficient for termwise Dirichlet subtraction.
Expand
\begin{equation}\label{hc:A}
 P(h_1,h_2,\ldots,h_R)=\sum_{j\ge0}A_j(\mathcal L)z^j,
             \qquad A_j\in\mathbb C[\mathcal L].
\end{equation}
Every fixed $A_j$ is a finite calculation. The degree in $\mathcal L$ is
bounded by $\deg_{X_1}P$.

\begin{proposition}[All principal coefficients]
\label{hc:principal}
At $s=1-j$, $j\ge0$, the entire Laurent principal part of $D_P$ is
\begin{equation}\label{hc:principal-formula}
 \sum_{\ell\ge0}
 \frac{\ell!\,[\mathcal L^\ell]A_j(\mathcal L)}{(s+j-1)^{\ell+1}}.
\end{equation}
Consequently $D_P$ is regular at all nonpositive even integers
exactly when $A_{2m+1}=0$ for every $m\ge0$.
\end{proposition}
\begin{proof}
Subtract finitely many terms of \eqref{hc:A} from the defining
series. Each subtracted term sums to
$(-1)^\ell\partial_s^\ell\zeta(s+j,1/2)$ times its coefficient.
The finite asymptotic remainder gives a normally convergent series
in a successively larger left half-plane. These formulas agree on
overlaps and continue $D_P$ meromorphically. The only principal
coefficient of $(-1)^\ell\partial_s^\ell\zeta(s+j,1/2)$ is
$\ell!/(s+j-1)^{\ell+1}$, proving the claim. Distinct logarithmic
degrees give distinct Laurent powers, so none is lost by merely
checking the residue.
\end{proof}

Here are the first three finite Bernoulli tests. Let $\partial_r$
denote polynomial differentiation with respect to $X_r$, omitting
indices larger than $R$, and define commuting operators
\begin{align*}
 A&=-\partial_2,&
 B&=\frac{\partial_1}{24}-\frac{\partial_3}{2},&
 C&=\frac{\partial_2}{12}-\frac{\partial_4}{3},\\
 D&=-\frac{7\partial_1}{960}+\frac{\partial_3}{8}
                    -\frac{\partial_5}{4},&
 E&=-\frac{7\partial_2}{240}+\frac{\partial_4}{6}
                    -\frac{\partial_6}{5}.
\end{align*}
Evaluate the following expressions at
$\mathbf b(\mathcal L)=(\mathcal L+\gamma,\zeta(2),\ldots,\zeta(R))$:
\begin{align}
 A_1(\mathcal L)&=(AP)(\mathbf b(\mathcal L)),\notag\\
 A_3(\mathcal L)&=\left((C+AB+A^3/6)P\right)(\mathbf b(\mathcal L)),\label{hc:first-tests}\\
 A_5(\mathcal L)&=\left((E+AD+BC+(A^2C+AB^2)/2
                   +A^3B/6+A^5/120)P\right)(\mathbf b(\mathcal L)).\notag
\end{align}
They follow by expanding the finite translation operator
$\exp(zA+z^2B+z^3C+z^4D+z^5E+\cdots)$.
For example,
\[
 A_3=\left(\frac{P_2}{12}-\frac{P_4}{3}
   -\frac{P_{12}}{24}+\frac{P_{23}}2-\frac{P_{222}}6\right)(\mathbf b(\mathcal L)).
\]
The rational-coefficient polynomial
$P=X_2^3-\tfrac{15}{2}X_2X_4$ satisfies $A_1=0$, because
$\zeta(4)=2\zeta(2)^2/5$, but
$A_3=\tfrac52\zeta(2)-1\ne0$. Thus its Dirichlet series is regular
at zero and has a simple pole at $-2$. More generally
$P=X_{2M+2}-\zeta(2M+2)$ passes all earlier tests and first has a
pole at $s=-2M$, with residue $-1/(2M+1)$. Finite lists of tested
spectral centers cannot replace the polynomial criterion when $R$
is unrestricted.

\subsection{Why the parity criterion is necessary}

The following elementary difference-algebra lemma is the relevant
form of the classical additive-difference mechanism behind
H\"older's theorem; compare the general theory in \cite{HardouinSinger}. It is proved here, so the formal-series application
does not rely on a converse assertion about analytic asymptotic
expansions.

\begin{lemma}[A polynomial relation forces rational telescoping]
\label{hc:additive}
Let $K=\mathbb C(x)$ with $\sigma x=x+1$, and let $E$ be a field
extension carrying an extension of $\sigma$. Suppose
$f_1,\ldots,f_d\in E$ satisfy $\sigma f_i=f_i+a_i$, $a_i\in K$.
If the $f_i$ are algebraically dependent over $K$, then there exist
constants $c_i\in\mathbb C$, not all zero, and $b\in K$ such that
\begin{equation}\label{hc:telescoping}
                       \sum_i c_i a_i=\sigma b-b.
\end{equation}
\end{lemma}
\begin{proof}
Choose a nonzero polynomial relation $Q(\mathbf f)=0$ of least
total degree $d_0$. Among such relations choose one with the fewest
nonzero monomials of degree $d_0$, and divide by one of its leading
coefficients. That coefficient is now one. The polynomial
\[
              \sigma Q(\mathbf X+\mathbf a)-Q(\mathbf X)
\]
also vanishes at $\mathbf f$. Its chosen leading monomial cancels,
and no new degree-$d_0$ monomial is created. Minimality first of
degree and then of leading support forces this difference to be
the zero polynomial. If $Q_{d_0}$ is the highest homogeneous part,
then $\sigma Q_{d_0}=Q_{d_0}$. Every rational periodic coefficient
is constant, so $Q_{d_0}\in\mathbb C[\mathbf X]$.
Comparison in degree $d_0-1$ gives
\[
 \sigma Q_{d_0-1}-Q_{d_0-1}
              =-\sum_i a_i\,\partial_i Q_{d_0}.
\]
At least one coefficient of the gradient on the right is a
nonzero vector in $\mathbb C^d$, since $d_0\ge1$ and the
characteristic is zero. Taking that monomial coefficient and
reversing the sign of the rational coefficient on the left proves
\eqref{hc:telescoping}.
\end{proof}

\begin{lemma}[Independence of the formal centered harmonic series]
\label{hc:independence}
For every finite $R$, the formal series $h_1,\ldots,h_R$ in
\eqref{hc:h1}--\eqref{hc:hr} are algebraically independent over
$\mathbb C(x)$, where $x=z^{-1}$ and $\mathcal L$ is a formal logarithm.
\end{lemma}
\begin{proof}
Work in $E=\mathbb C((z))(\mathcal L)$. Extend translation by
\[
 \sigma z=\frac{z}{1+z},\qquad
 \sigma \mathcal L=\mathcal L+\log(1+z),
\]
and use the derivation $\mathcal D z=-z^2$, $\mathcal D \mathcal L=z$.
These operations commute. Put $f=h_1-\gamma$.
The Bernoulli expansion and the digamma recurrence give the formal
identity
\[
 \sigma f-f=\frac1{x+1/2},\qquad
 \sigma\mathcal D^j f-\mathcal D^j f
       =\frac{(-1)^j j!}{(x+1/2)^{j+1}}\quad(j\ge0).
\]
It can equally be checked coefficientwise from the Bernoulli
generating function. A nonzero constant linear combination of the
right sides has exactly one finite pole, at $-1/2$.

No nonzero rational difference $b(x+1)-b(x)$ has just one finite
pole. Indeed partition the finite poles of $b$ into classes modulo
$\mathbb Z$. In any nonempty class, the extreme poles of $b$ give
two distinct, uncanceled extreme poles of its difference. If $b$
has no finite poles, it is polynomial and its difference is
polynomial. Lemma~\ref{hc:additive} therefore proves algebraic
independence of $f,\mathcal Df,\ldots,\mathcal D^{R-1}f$.
Finally
\[
 h_1=f+\gamma,\qquad
 h_r=\zeta(r)-\frac{(-1)^r}{(r-1)!}\mathcal D^{r-1}f\quad(r\ge2)
\]
is an invertible affine change of these variables.
\end{proof}

\begin{proof}[Proof of Theorem~\ref{hc:main}]
Let $\tau$ act on formal expansions by $z\mapsto-z$, $\mathcal L\mapsto \mathcal L$.
Equations \eqref{hc:h1}--\eqref{hc:hr} show
$\tau h_r=\iota(h_r)$: the even harmonic tails change sign and
the odd harmonic variables do not. By
Proposition~\ref{hc:principal}, absence of every nonpositive even
principal part is equivalent to
$\tau P(\mathbf h)=P(\mathbf h)$, hence to
$(\iota P-P)(\mathbf h)=0$. Lemma~\ref{hc:independence} makes
this equivalent to the polynomial identity $\iota P=P$.
In centered variables $Y_{2k}$ the involution changes all signs
simultaneously. Its fixed monomials have even total $Y$-degree,
and each such monomial factors into quadratic pairs. This proves
\eqref{hc:invariant-ring}.
\end{proof}

\begin{corollary}[Algebraic coefficients]
\label{hc:algebraic}
If $P$ has algebraic-number coefficients, the criterion holds if
and only if $P$ is independent of every even-index variable.
\end{corollary}
\begin{proof}
Write $\zeta(2k)=c_k\pi^{2k}$ with nonzero $c_k\in\mathbb Q$.
Since $\pi^2$ is transcendental, invariance of a polynomial over
$\overline{\mathbb Q}$ under \eqref{hc:involution} implies its
invariance under the polynomial reflection
$X_{2k}\mapsto2c_kT^k-X_{2k}$ for every indeterminate $T$.
Composing two such reflections gives invariance under translations
by $2(c_kT^k-c_kU^k)_k$. Differentiating in $T$ and then setting
$T=U$ yields
\[
             \sum_k k c_k U^{k-1}\partial_{2k}P=0.
\]
Distinct powers of $U$ force each $\partial_{2k}P$ to vanish.
The converse follows directly from Theorem~\ref{hc:main}.
\end{proof}

This corollary uses only the established transcendence of $\pi$.
The main theorem makes no assertion about algebraic independence
of evaluated zeta or Stieltjes constants.

\subsection{All centered even moments of a trigamma square}

The invariant numerator $(H_n^{(2)}-\zeta(2))^2$ admits a particularly
short all-degree value formula. Put
\[
 T_n=\zeta(2,n+1)=\psi'(n+1),\qquad
 V_M=D_{(X_2-\zeta(2))^2}(-2M),\quad M\ge0.
\]
Every point in this definition is regular by Theorem~\ref{hc:main}.

\begin{theorem}[Finite Bernoulli formula for all even moments]
\label{hc:trigamma}
For every $M\ge0$,
\begin{align}\label{hc:all-moments}
 V_M={}&3B_{2M}(1/2)\zeta(3)\notag\\
 &-\frac1{2(2M+1)}\sum_{\ell=1}^{M}
 \binom{2M+1}{2M-2\ell}B_{2M-2\ell}(1/2)
                     (2\ell-1)B_{2\ell-2}.
\end{align}
The sum is empty at $M=0$. In particular all these values belong
to $\mathbb Q+\mathbb Q\zeta(3)$.
\end{theorem}
\begin{proof}
The polynomial partial sum is
\[
 S_M(k)=\sum_{n=0}^{k-1}(n+1/2)^{2M}
 =\frac{B_{2M+1}(k+1/2)}{2M+1}
 =\sum_{\ell=0}^M s_{M,\ell} k^{2\ell+1},
\]
where
$s_{M,\ell}=\binom{2M+1}{2M-2\ell}B_{2M-2\ell}(1/2)/(2M+1)$.
Since $T_{k-1}-T_k=k^{-2}$, finite summation by parts gives
\[
 \sum_{n=0}^{N-1}(n+1/2)^{2M}T_n^2
 =\sum_{\ell=0}^M s_{M,\ell}E_{2\ell+1}(N),
\]
with
\[
 E_q(N)=N^qT_N^2+2\sum_{k=1}^N k^{q-2}T_{k-1}
                                -\sum_{k=1}^N k^{q-4}.
\]
At $q=1$ the limit is $3\zeta(3)$: the classical Euler sum \cite{Can}
$\sum_{k\ge1}H_k/k^2=2\zeta(3)$ implies
$\sum_{k\ge1}T_{k-1}/k=2\zeta(3)$, by interchanging a positive
double sum. The remaining diagonal term is $\zeta(3)$.

For odd $q\ge3$, let $P_{q-2}(N)=\sum_{k=1}^N k^{q-2}$.
Interchanging a finite sum gives
\[
 \sum_{k=1}^N k^{q-2}T_{k-1}
   =\sum_{j=1}^N\frac{P_{q-2}(j)}{j^2}
                                +P_{q-2}(N)T_N.
\]
In Faulhaber's polynomial the half-leading term cancels
$\sum k^{q-4}$ from $E_q$. All other finite power sums are
polynomials in $N$ with zero constant coefficient. Hence the
constant term of the expansion of $E_q(N)$ at infinity is the
constant of $N^qT_N^2+2P_{q-2}(N)T_N$.
Use
\[
 T_N=N^{-1}-\tfrac12N^{-2}
                   +\sum_{j\ge1}B_{2j}N^{-2j-1}.
\]
Only finitely many terms are needed. The first product contributes
$-B_{q-3}$ and the second contributes
$B_{q-3}-(q-2)B_{q-3}/2$. Their sum is
$-(q-2)B_{q-3}/2$.

It remains to identify this cutoff constant with the stated
spectral value. The centered expansion of $T_n^2$ contains only
even powers of $(n+1/2)^{-1}$. Subtract from
$(n+1/2)^{2M}T_n^2$ all its nondecaying even powers. The residual
is $O(n^{-2})$ and is ordinarily summable. Each subtracted power
has spectral value $\zeta(-2j,1/2)=0$, and its finite partial sum
is an odd polynomial in $N$ with zero constant. Thus the residual
sum, the cutoff constant in $N$, and $V_M$ coincide. Combining
the coefficients $s_{M,\ell}$ proves \eqref{hc:all-moments}.
\end{proof}

The first four identities can be written as ordinary convergent
bracketed series, with $x_n=n+1/2$:
\begin{align}
 \sum_{n\ge0}T_n^2&=3\zeta(3),\notag\\
 \sum_{n\ge0}\{x_n^2T_n^2-1\}
        &=-\frac16-\frac14\zeta(3),\notag\\
 \sum_{n\ge0}\{x_n^4T_n^2-x_n^2+\tfrac16\}
        &=\frac1{30}+\frac7{80}\zeta(3),\label{hc:moment-examples}\\
 \sum_{n\ge0}\{x_n^6T_n^2-x_n^4+\tfrac16x_n^2-\tfrac{47}{720}\}
        &=\frac1{672}-\frac{31}{448}\zeta(3).\notag
\end{align}
Each bracket is $O(n^{-2})$. These statements do not assign an
ordinary sum to any of the divergent unbracketed expressions.



% Begin sections/04_collision
\section{Geometric collisions at Stieltjes index zero}
\label{gc:section}

The periodic contact reports \cite{RepoPeriodic,RepoContact,RepoMixed} develop
two-factor and spectral completions. Here the shifts themselves move toward
collision. The result covers every finite number of factors and every
argument-derivative order at Stieltjes index zero.

Set $f(x)=\gamma_0(x)=-\psi(x)$ and $h(x)=f(1+x)$.  The Hurwitz shift
identity gives the exact local decomposition
\begin{equation}\label{gc:one-sided}
 f(\{x\})=\frac{\boldsymbol{1}_{x>0}}x+h(x),\qquad -1<x<1,
 \quad x\ne0.
\end{equation}
Here the negative half of the formula uses $\{x\}=1+x$.  The germ $h$ is
holomorphic on $|x|<1$, with
\begin{equation}\label{gc:hseries}
 h(x)=\gamma+\sum_{n\ge1}(-1)^n\zeta(n+1)x^n.
\end{equation}
The finite-part convention is fixed throughout: at each right-hand
singular point $a$ remove $(a,a+\eta)$ and take the constant term in the
\emph{ambient coordinate} $\eta$.  Independent cutoffs are used at distinct
points.  No collision-dependent rescaling of these coordinates is made.

For distinct sufficiently small real shifts $s_1,\ldots,s_d$, define
\begin{equation}\label{gc:Cdef}
 C_d(\boldsymbol{s})=
 \operatorname{FP}\int_{\mathbb R/\mathbb Z}
       \prod_{i=1}^d f(\{x+s_i\})\,dx,
 \qquad
 Q_d=\operatorname{FP}\int_0^1 f(x)^d\,dx.
\end{equation}
For a nonempty subset $S\subseteq\{1,\ldots,d\}$, let $m(S)$ denote its
index with smallest shift.  It is a fixed index when an ordering chamber
is fixed.

\begin{theorem}[Complete local counterterm for any number of factors]
\label{gc:main}
Define
\begin{equation}\label{gc:T}
 T_d(\boldsymbol{s})=
 -\sum_{\substack{S\subseteq\{1,\ldots,d\}\\ |S|\ge2}}
  \ \sum_{i\in S\setminus\{m(S)\}}
  \frac{\displaystyle\prod_{j\notin S}h(s_j-s_i)}
       {\displaystyle\prod_{k\in S\setminus\{i\}}(s_k-s_i)}
       \log(s_i-s_{m(S)}).
\end{equation}
On each fixed ordering chamber,
$C_d(\boldsymbol{s})-T_d(\boldsymbol{s})$ is the restriction of a real
analytic function on a full neighborhood of $\boldsymbol{s}=0$.  Its value
at total collision is $Q_d$.  Consequently
\begin{equation}\label{gc:anypath}
 C_d(\boldsymbol{s})-T_d(\boldsymbol{s})\longrightarrow Q_d
 \qquad (\boldsymbol{s}\longrightarrow0)
\end{equation}
along every path remaining in that chamber, without a lower bound on the
ratios of the gaps.
\end{theorem}

\begin{proof}
Choose $\delta>0$ small enough that all required germs are holomorphic on
$|z|<3\delta$, and restrict the shifts to $|s_i|<\delta/8$.  Use the circle
coordinate $-\delta<x<1-\delta$.  Outside $(-\delta,\delta)$ the entire
integrand is jointly analytic in the shifts, and its integral is analytic.
On $(-\delta,\delta)$, expand the exact product of
\eqref{gc:one-sided} by the subset $S$ of factors contributing their pole.
The empty subset is analytic.  For a nonempty $S$, the product of the
indicator functions is $\boldsymbol{1}_{x>-s_{m(S)}}$, and the contribution
is the finite part of
\begin{equation}\label{gc:subset-integral}
 \int_{-s_{m(S)}}^\delta \frac{q_S(x;\boldsymbol{s})}
                                      {D_S(x;\boldsymbol{s})}\,dx,
 \quad
 q_S(x;\boldsymbol{s})=\prod_{j\notin S}h(x+s_j),
 \quad D_S(x;\boldsymbol{s})=\prod_{i\in S}(x+s_i).
\end{equation}

Let $P_S$ be the interpolation polynomial of degree less than $|S|$ which
agrees with $q_S$ at the nodes $-s_i$, $i\in S$.  Both $P_S$ and
$R_S=(q_S-P_S)/D_S$ extend holomorphically when nodes coalesce.  One
explicit justification is the fixed-contour expression
\begin{equation}\label{gc:hermite}
 P_S(x;\boldsymbol{s})=\frac1{2\pi i}\oint_{|z|=2\delta}
 \frac{q_S(z;\boldsymbol{s})}{D_S(z;\boldsymbol{s})}
 \frac{D_S(z;\boldsymbol{s})-D_S(x;\boldsymbol{s})}{z-x}\,dz.
\end{equation}
Its integrand is holomorphic in the parameters and the quotient of
polynomials is a polynomial in $x$ of the required degree.  Cauchy's
formula gives the jointly holomorphic remainder
\[
 R_S(x;\boldsymbol{s})=\frac1{2\pi i}\oint_{|z|=2\delta}
 \frac{q_S(z;\boldsymbol{s})}
 {D_S(z;\boldsymbol{s})(z-x)}\,dz,
 \qquad |x|<2\delta.
\]

For distinct nodes the rational principal part is
\[
 \frac{P_S(x;\boldsymbol{s})}{D_S(x;\boldsymbol{s})}
 =\sum_{i\in S}\frac{b_{S,i}(\boldsymbol{s})}{x+s_i},
 \qquad
 b_{S,i}(\boldsymbol{s})=
 \frac{q_S(-s_i;\boldsymbol{s})}
      {\prod_{k\in S\setminus\{i\}}(s_k-s_i)}.
\]
Its finite-part integral in \eqref{gc:subset-integral} equals
\begin{equation}\label{gc:integrated}
 \sum_{i\in S}b_{S,i}\log(\delta+s_i)
 -\sum_{i\in S\setminus\{m(S)\}}
       b_{S,i}\log(s_i-s_{m(S)}).
\end{equation}
At the minimal node the discarded primitive is $b_{S,m(S)}\log\eta$,
whose constant term is zero under the specified coordinate convention.
The second sum in \eqref{gc:integrated} is precisely the corresponding
summand of \eqref{gc:T}.  Although the individual coefficients $b_{S,i}$
have poles at coincident nodes, the first sum has the holomorphic form
\begin{equation}\label{gc:upper-contour}
 \frac1{2\pi i}\oint_{|z|=\delta/2}
 \frac{q_S(z;\boldsymbol{s})}{D_S(z;\boldsymbol{s})}
                      \operatorname{Log}(\delta-z)\,dz.
\end{equation}
The logarithm is the holomorphic branch equal to $\log\delta$ at $z=0$.
This contour is deliberately smaller than the interpolation contour in
\eqref{gc:hermite}.  The integral of $R_S$ from $-s_{m(S)}$ to $\delta$ is
also holomorphic, since its lower endpoint is a fixed coordinate function
on the chosen chamber.

At $\boldsymbol{s}=0$, $D_S=x^{|S|}$ and $P_S$ is the Taylor polynomial of
$q_S$ through degree $|S|-1$.  Formula \eqref{gc:upper-contour} is exactly
the finite-part integral of $P_S/x^{|S|}$ from $0$ to $\delta$; the
remainder gives the ordinary integral of the analytic part.  Summing the
subsets therefore gives the same local finite part as the directly
merged product.  The outer integral supplies its remaining part, proving
that the analytic remainder has value $Q_d$.
\end{proof}

\begin{remark}[Meaning of the chamber qualification]
The extension asserted above is an extension \emph{from each chamber}.
Different chambers choose different lower-endpoint indices in
\eqref{gc:subset-integral}; no assertion that the resulting analytic germs
agree off the total diagonal is needed.  All their values at total
collision are the prescribed merged finite part.  In particular the
theorem proves a path-independent limit only after subtracting the full
$T_d$, including any finite constant it contributes on the chosen path.
\end{remark}

\subsection{Every argument-derivative order}
For nonnegative integers $r_1,\ldots,r_d$, put
\[
 d_i=r_i+1,\qquad A_i=(-1)^{r_i}r_i!,\qquad
 h_i(x)=\gamma_0^{(r_i)}(1+x).
\]
Ordinary differentiation away from the integer endpoints gives
\[
 \gamma_0^{(r_i)}(\{x\})
    =A_i\boldsymbol{1}_{x>0}x^{-d_i}+h_i(x).
\]
Define $C_{\boldsymbol r}(\boldsymbol s)$ and $Q_{\boldsymbol r}$ by
replacing the factors in \eqref{gc:Cdef} by the indicated argument
derivatives.  For a nonempty $S$, set
\[
 A_S=\prod_{i\in S}A_i,\qquad
 q_S(x;\boldsymbol s)=\prod_{j\notin S}h_j(x+s_j),\qquad
 D_S(x;\boldsymbol s)=\prod_{i\in S}(x+s_i)^{d_i}.
\]
For $i\in S$ and $1\le q\le d_i$, let
\begin{equation}\label{gc:B}
 B_{S,i,q}(\boldsymbol s)=
 \frac1{(d_i-q)!}
 \left.\left(\frac{d}{dx}\right)^{d_i-q}
 \frac{q_S(x;\boldsymbol s)}
 {\prod_{k\in S\setminus\{i\}}(x+s_k)^{d_k}}
 \right|_{x=-s_i}.
\end{equation}

\begin{theorem}[Repeated-pole collision formula]\label{gc:derivative}
Let
\begin{align}\label{gc:Tr}
 T_{\boldsymbol r}(\boldsymbol s)
 ={}&\sum_{\substack{S\subseteq\{1,\ldots,d\}\\ |S|\ge2}}
 A_S\sum_{i\in S\setminus\{m(S)\}}
 \left\{-B_{S,i,1}\log(s_i-s_{m(S)})\right.\notag\\[-2pt]
 &\hspace{42mm}\left.
 +\sum_{q=2}^{d_i}\frac{B_{S,i,q}}{q-1}
                   (s_i-s_{m(S)})^{1-q}\right\}.
\end{align}
Then $C_{\boldsymbol r}-T_{\boldsymbol r}$ has a real analytic extension
from each ordering chamber through total collision, with value
$Q_{\boldsymbol r}$ there.
\end{theorem}

\begin{proof}
Repeat the preceding proof with Hermite interpolation of multiplicity
$d_i$ at $-s_i$.  Formula \eqref{gc:hermite}, with the new $D_S$, proves
holomorphy of the interpolation polynomial and its remainder, including
coalescing nodes.  The principal part is
$\sum_{i\in S}\sum_{q=1}^{d_i}B_{S,i,q}(x+s_i)^{-q}$.  Its upper primitive
is again \eqref{gc:upper-contour}, multiplied by $A_S$, since
\[
 \frac1{(q-1)!}\frac{d^{q-1}}{dz^{q-1}}
      \operatorname{Log}(\delta-z)
 =-\frac{(\delta-z)^{1-q}}{q-1},\qquad q\ge2.
\]
At a nonminimal node, subtracting the lower primitive gives exactly the
bracket in \eqref{gc:Tr}.  At the minimal node each lower primitive is a
pure power of $\eta$ or a multiple of $\log\eta$, and hence has zero
constant term.  The same coalescence argument identifies the value of the
analytic remainder with $Q_{\boldsymbol r}$.
\end{proof}

This proof differentiates only ordinary local analytic functions in
\eqref{gc:B}.  Differentiating an already extended periodic distribution
can generate contact terms, and does not supply a substitute for the
direct repeated-pole proof.  For example, with two factors and shifts
$(0,a)$, the formula gives
\[
 T_{(1,0)}(0,a)=\frac{\log a}{a^2},\qquad
 T_{(0,1)}(0,a)=\frac{1-\log a}{a^2}.
\]
Both analytic remainders tend to
$Q_{(1,0)}=Q_{(0,1)}=\zeta(2)$, which also follows by integrating
$ff'=(f^2)'/2$ and taking the endpoint finite part.
The separated integrals satisfy the exact identity
\begin{equation}\label{gc:pair-sum}
 C_{(1,0)}(0,a)+C_{(0,1)}(0,a)
   =\frac{\pi^2}{\sin^2(\pi a)},\qquad 0<a<1.
\end{equation}
Indeed, apply the fundamental theorem of calculus to the product
$f(x)f(\{x+a\})$ on the two nonsingular intervals.  At $0$ the left
boundary value minus the constant in the right Laurent expansion is
$-f'(a)$; at $1-a$ it is $-f'(1-a)$.  Taking the finite parts and using
trigamma reflection gives
$-f'(a)-f'(1-a)=\psi'(a)+\psi'(1-a)=\pi^2/\sin^2(\pi a)$.
Thus \eqref{gc:pair-sum} records the boundary correction required by the
specified finite part.  The sum of the two local counterterms is $a^{-2}$,
and the sum of their limiting remainders is $2\zeta(2)=\pi^2/3$, in
agreement with the Taylor expansion of its right side.

\subsection{Fixed rates and a finite zeta--logarithm expression}
For $s_i=a_i\epsilon$ with fixed distinct real rates, define
$m(S)=\operatorname*{argmin}_{i\in S}a_i$.  Coefficient extraction from
\eqref{gc:T} gives
\begin{align}\label{gc:fixed-constant}
 \operatorname{CT}_{\epsilon\downarrow0}C_d(\epsilon\boldsymbol a)
 ={}&Q_d-\sum_{\substack{S\subseteq\{1,\ldots,d\}\\ |S|\ge2}}
 \ \sum_{i\in S\setminus\{m(S)\}}
 \frac{\log(a_i-a_{m(S)})}
      {\prod_{k\in S\setminus\{i\}}(a_k-a_i)}\notag\\
 &\hspace{18mm}\cdot
 [\epsilon^{|S|-1}]
       \prod_{j\notin S}h\bigl(\epsilon(a_j-a_i)\bigr).
\end{align}
Here the constant-term operation discards positive powers of
$\log\epsilon$ even when their power of $\epsilon$ is zero.  Only finitely
many Taylor coefficients of \eqref{gc:hseries} occur.  In particular the
finite anomaly is an explicit polynomial in Euler's constant and ordinary
zeta values, with rational functions of the rates and logarithms of their
positive differences as coefficients.  This is a local conclusion; it
does not evaluate the merged global constant $Q_d$.

For equally spaced triples, \eqref{gc:fixed-constant} recovers the earlier
three-factor contact specialization \cite{RepoMixed}:
\[
 \operatorname{CT}_{\epsilon\downarrow0}
 C_3(0,\epsilon,2\epsilon)=Q_3+\frac{\zeta(2)}2\log2.
\]
For four equally spaced points, the new formula is
\begin{equation}\label{gc:quartic-example}
 \operatorname{CT}_{\epsilon\downarrow0}
 C_4(0,\epsilon,2\epsilon,3\epsilon)
 =Q_4+\gamma\zeta(2)(2\log2+\log3)
       -\frac32\zeta(3)(3\log2+\log3).
\end{equation}

\subsection{A hierarchical collision}
Take shifts $(0,\epsilon,\epsilon+\epsilon^2)$, so the last pair approaches
one another on a smaller scale than its distance from the first point.
The four terms in \eqref{gc:T} are exactly
\begin{align*}
 T_{\{1,2,3\}}&=\frac{\log\epsilon}{\epsilon^3}
 -\frac{\log(\epsilon+\epsilon^2)}{\epsilon^3(1+\epsilon)},\\
 T_{\{1,2\}}&=\frac{h(\epsilon^2)}\epsilon\log\epsilon,\\
 T_{\{1,3\}}&=\frac{h(-\epsilon^2)}{\epsilon+\epsilon^2}
                         \log(\epsilon+\epsilon^2),\\
 T_{\{2,3\}}&=\frac{h(-\epsilon-\epsilon^2)}{\epsilon^2}
                         \log(\epsilon^2).
\end{align*}
Writing $\Lambda=\log\epsilon$, expanding these exact terms, and applying
Theorem~\ref{gc:main} yields
\begin{align}\label{gc:hierarchy}
 C_3(0,\epsilon,\epsilon+\epsilon^2)
 ={}&\frac{(1+2\gamma)\Lambda-1}{\epsilon^2}
 +\frac{(2\gamma+2\zeta(2)-1)\Lambda+\tfrac32}{\epsilon}\notag\\
 &+(1-\gamma+2\zeta(2)+2\zeta(3))\Lambda
 +Q_3+\gamma-\frac{11}{6}
 +O(\epsilon|\log\epsilon|).
\end{align}
The finite constant is therefore $Q_3+\gamma-11/6$.  The rational term
$-11/6$ comes from the all-pole subset, while the term $\gamma$ comes from
$T_{\{1,3\}}$.  This example shows why one cannot obtain hierarchical
constant terms by retaining only the singular powers of a fixed-rate
expansion and subsequently varying its rate parameters.

\subsection{Verification and precise scope}
The companion script \texttt{verify\_collision.py} verifies the quartic
constant and the full hierarchical expansion through order zero by exact
symbolic algebra.  It also evaluates the separated circle integral by
subtracting its actual right-hand poles interval by interval, independently
of \eqref{gc:T}.  Direct Laurent subtraction at
the merged endpoint gives
\begin{align*}
 Q_3&=-1.9662627343915343406643753384116\ldots,\\
 Q_4&=\phantom{-}13.510120376015827411606240712830\ldots.
\end{align*}
The computed differences $C_d-T_d-Q_d$ tend to zero linearly on both
displayed paths, at four values of $\epsilon$ for each path and $55$
working decimal digits.  The separate script
\texttt{verify\_derivative\_pair.py} checks both displayed derivative
counterterms at the same four gaps, using independent double-pole
subtraction and $65$ working decimal digits.  Its four comparisons with
\eqref{gc:pair-sum} have observed absolute residual at most
$6.4\times10^{-51}$.  The package includes the raw JSON records.
These floating-point diagnostics check implementation and signs; the
interpolation proof supplies the theorem.

The result settles the arbitrary-factor and arbitrary argument-derivative
parts of the unit-frequency geometric collision question at Stieltjes index zero,
including hierarchical approaches.  Higher Stieltjes indices introduce
logarithmic singular germs and are not covered by the rational
interpolation argument.  Associativity of successively renormalized
partial mergers also remains a separate question: an exact simultaneous
subtraction does not by itself identify every iterated subtraction rule.
Two further extensions are a complete statement for integer-frequency
singular grids, with each local pole amplitude scaled in the ambient
coordinate, and a weighted version for analytic periodic test functions
which would provide collision formulas for Fourier moments.  Neither is
included in the theorems proved here.


% Begin sections/05_tornheim
\section{The complete quartic directional layer of Tornheim zeta}
\label{qt:sec}

This section answers the explicit quartic-layer question in the incoming
\emph{Cubic Tornheim Derivatives and Harmonic Stieltjes Constants},
Section 7, Question 4 \cite{RepoCubic}, and the corresponding question in
\emph{Independent Orders and Exact Reductions}, Section 7 \cite{RepoIndependent}.
The incoming work already proves the cubic ray formula and evaluates its
diagonal coordinate in a convergent harmonic series. Those results are
not new here. The present extension classifies the next homogeneous layer,
gives convergent representatives for every coordinate, and obtains exact
quartic cancellation identities.

Put
\[
 L=\log(2\pi),\qquad d_j=\zeta^{(j)}(0),\qquad
 e_j=\zeta^{(j)}(-1),\qquad Z_j=\zeta(j),\qquad g=\gamma.
\]
Spectral derivatives are denoted by parentheses on zeta, whereas
subscripts on a multivariable analytic function mean partial derivatives.
The classical Mellin continuation \cite{BorweinDilcher} of
\[
 T(A,B,C)=\sum_{m,n\geq1}m^{-A}n^{-B}(m+n)^{-C}
\]
has the following local form, already established in the incoming
independent-order report \cite{RepoIndependent}:
\begin{equation}\label{qt:polar}
 T(A,B,C)=\zeta(A)\zeta(B)
 -\frac{C\zeta(B-1)}{A+C}
 -\frac{C\zeta(A-1)}{B+C}+C R(A,B,C),
\end{equation}
where $R$ is holomorphic near the origin and symmetric in $A,B$.
The two indicated polar quotients must be restricted before evaluating
at the origin. Thus for $(a+c)(b+c)\ne0$ the function
$W_{a,b,c}(t)=T(at,bt,ct)$ is holomorphic near zero, although $T$ is not
jointly holomorphic there.

\subsection{Six cubic coefficients and three exact constraints}

All derivatives in the next display are evaluated at $(0,0,0)$:
\[
 \alpha=R_{AAA},\quad \beta=R_{AAB},\quad
 \delta=R_{AAC},\quad \varepsilon=R_{ABC},\quad
 \varphi=R_{ACC},\quad \chi=R_{CCC}.
\]
Symmetry gives the corresponding derivatives with $A,B$ interchanged.
These six quantities describe the cubic homogeneous part of $R$.

\begin{proposition}\label{qt:constraints}
One has
\begin{equation}\label{qt:constraints-eq}
 \chi=\frac{e_4-d_4}{4},\qquad
 \delta=\frac{d_2^2-d_4}{4},\qquad
 \alpha+3\varphi=-\frac{Ld_3}{2}-d_4.
\end{equation}
If $\Omega_4=W_{1,1,1}^{(4)}(0)$, then
\begin{equation}\label{qt:diagonal}
 \boxed{\Omega_4=24(\beta+\varepsilon)-8Ld_3+12d_2^2-16d_4.}
\end{equation}
\end{proposition}
\begin{proof}
The exact axis $T(0,0,C)=\zeta(C-1)-\zeta(C)$ gives
\[
 C R(0,0,C)=\zeta(C-1)-\zeta(C)-\frac5{12}.
\]
Its coefficient of $C^4$ proves the first assertion.
The Euler stuffle identity on the plane $B=0$ gives
\begin{align}\label{qt:stuffle}
 C R(A,0,C)+A R(C,0,A)
 ={}&\zeta(A)\zeta(C)-\zeta(A+C)
 +\tfrac12\bigl(\zeta(A)+\zeta(C)\bigr)\notag\\
 &+\zeta(-1)+\zeta(A-1)+\zeta(C-1).
\end{align}
The coefficient of $A^2C^2$ is $\delta$ on the left and
$(d_2^2-d_4)/4$ on the right. The coefficient of $A^3C$ is
$\alpha/6+\varphi/2$ on the left and
$d_1d_3/6-d_4/6$ on the right. Since $d_1=-L/2$, this proves
the remaining two constraints.

On the diagonal, \eqref{qt:polar} becomes
$W_{1,1,1}(t)=\zeta(t)^2-\zeta(t-1)+tR(t,t,t)$.
Its fourth derivative is
\[
 -d_4-4Ld_3+6d_2^2-e_4
 +8\alpha+24\beta+24\delta+24\varepsilon+24\varphi+4\chi.
\]
Substituting \eqref{qt:constraints-eq} proves \eqref{qt:diagonal}.
\end{proof}

The three constraints have rank three in the six-dimensional space of
formal cubic jets symmetric in $A,B$. Thus that formal reconstruction
problem has exactly three free coordinates. This dimension statement
does not imply arithmetic independence of the evaluated constants.

\subsection{An explicit formula for every fourth ray derivative}

\begin{theorem}\label{qt:main}
For $a,b,c\in\mathbb C$ with $(a+c)(b+c)\ne0$,
\begin{align}\label{qt:all-rays}
 W_{a,b,c}^{(4)}(0)={}&abc^2\Omega_4
 +4c(a+b)(a^2-ab+b^2-c^2)\alpha\notag\\
 &+12abc(a+b-2c)\beta\notag\\
 &-2\bigl[ab(a^2+b^2)-4abc^2+c^3(a+b)\bigr]Ld_3\notag\\
 &+3\bigl[2a^2b^2+c^2(a^2+b^2-4ab)\bigr]d_2^2\notag\\
 &+\bigl[-\tfrac12(a^4+b^4)-3c^2(a^2+b^2)
              +16abc^2-4c^3(a+b)-c^4\bigr]d_4\notag\\
 &+\left[c^4-\frac{cb^4}{a+c}-\frac{ca^4}{b+c}\right]e_4.
\end{align}
The coordinates $\alpha,\beta,\Omega_4$ have the convergent
Gamma and polylogarithmic representations below; no unspecified partial
derivative is required to evaluate the formula.
\end{theorem}
\begin{proof}
Expanding the holomorphic restriction of \eqref{qt:polar} gives
\begin{align*}
 W_{a,b,c}^{(4)}(0)={}&-\tfrac12(a^4+b^4)d_4
 -2ab(a^2+b^2)Ld_3+6a^2b^2d_2^2\\
 &-c\left(\frac{b^4}{a+c}+\frac{a^4}{b+c}\right)e_4
 +4c\alpha(a^3+b^3)+12c\beta ab(a+b)\\
 &+12c^2\delta(a^2+b^2)+24abc^2\varepsilon
 +12c^3\varphi(a+b)+4c^4\chi.
\end{align*}
Eliminate $\delta,\varphi,\chi$ using \eqref{qt:constraints-eq},
then eliminate $\varepsilon$ with \eqref{qt:diagonal}. Factoring the
$\alpha$ coefficient proves the claimed formula.
\end{proof}

For instance,
\begin{align}
 W_{1,1,2}^{(4)}(0)={}&4\Omega_4-48\alpha-48\beta-4Ld_3
                      -18d_2^2-41d_4+\tfrac{44}{3}e_4,
 \label{qt:112}\\
 W_{1,2,2}^{(4)}(0)={}&8\Omega_4-24\alpha-48\beta-4Ld_3
                      -12d_2^2-\tfrac{105}{2}d_4+\tfrac{29}{6}e_4.
 \label{qt:122}
\end{align}
The three positive anchor rays $(1,1,1),(1,1,2),(1,2,2)$ recover
all three formal coordinates: the determinant of their coefficient
matrix in $(\Omega_4,\alpha,\beta)$ is $1152$.

\subsection{Convergent representatives: a Gamma moment and two kernels}

The incoming boundary-jet hierarchy \cite{RepoIndependent} already gives, with
\[
 \rho(n)=\log\Gamma(n+1)-(n+\tfrac12)\log n+n-\frac L2-\frac1{12n},
\]
the identity
\begin{equation}\label{qt:alpha}
 \alpha=-\sum_{n\ge1}\rho(n)\log^3n-e_4-\tfrac12d_4-e_3-\frac{\gamma_3}{12}.
\end{equation}
The series converges absolutely because $\rho(n)=O(n^{-3})$.
Its use as a quartic coordinate is new here; the all-order boundary
identity itself is credited to the incoming report.

For the other two coordinates, define for $x>0$
\[
 f_j(x)=\left.\partial_s^j\operatorname{Li}_s(e^{-x})\right|_{s=0},
 \qquad
 p_1(x)=\frac{g+\log x}{x},\quad
 p_2(x)=\frac{(g+\log x)^2+Z_2}{x}.
\]
Put
\begin{align*}
 S_{21}(x)&=p_2(x)p_1(x)+p_2(x)d_1+p_1(x)d_2
                  -x\bigl(p_2(x)e_1+p_1(x)e_2\bigr)+d_2d_1,\\
 S_{11}(x)&=p_1(x)^2+2p_1(x)d_1-2xp_1(x)e_1+d_1^2.
\end{align*}
The ordinary convergent integrals are
\begin{align}\label{qt:J21}
 J_{21}={}&\int_0^1\frac{f_2(x)f_1(x)-S_{21}(x)}x\,dx
                  +\int_1^\infty\frac{f_2(x)f_1(x)}x\,dx,\\
 K_{11}={}&\int_0^1\frac{\log x+g}{x}
                    \bigl(f_1(x)^2-S_{11}(x)\bigr)\,dx
                  +\int_1^\infty\frac{\log x+g}{x}f_1(x)^2\,dx.
 \label{qt:K11}
\end{align}
Set $\mu_2=g^2+Z_2$ and $\mu_3=g^3+3gZ_2+2Z_3$, and define
\begin{align}\label{qt:constants}
 E_{21}={}&gd_1d_2-(\mu_2+2g+2)d_1-(g+1)d_2
                -\frac{\mu_3e_1}{3}-\frac{\mu_2e_2}{2}\notag\\
 &-\frac38-\frac{3g}{4}-\frac{\mu_2}{4}
                -\frac{g^2}{2}-\frac{g\mu_2}{2},\\
 E_{111}={}&-\frac{g^3}{2}-\frac{3g^2}{4}-\frac{3g}{4}-\frac38
               -2(g^2+2g+2)d_1\notag\\
 &+\frac{g^2-Z_2}{2}d_1^2-\frac{2(g^3-Z_3)}{3}e_1.
\end{align}

\begin{proposition}\label{qt:integrals}
The integral and constant formulas give
\begin{equation}\label{qt:beta-epsilon}
 \boxed{\beta=J_{21}+E_{21},\qquad
        \varepsilon=K_{11}+E_{111}.}
\end{equation}
Together with \eqref{qt:diagonal}, these evaluate $\Omega_4$ using
two ordinary, absolutely convergent one-dimensional integrals.
\end{proposition}
\begin{proof}
For parameters near zero put $G(C)=1/\Gamma(1+C)$ and
\begin{align*}
 S(A,B;x)={}&\Gamma(1-A)\Gamma(1-B)x^{A+B-2}
 +\Gamma(1-A)\zeta(B)x^{A-1}
 +\Gamma(1-B)\zeta(A)x^{B-1}\\
 &-\Gamma(1-A)\zeta(B-1)x^A
 -\Gamma(1-B)\zeta(A-1)x^B+\zeta(A)\zeta(B),\\
 I(A,B,C)={}&\int_0^1x^{C-1}
  \bigl(\operatorname{Li}_A(e^{-x})\operatorname{Li}_B(e^{-x})-S(A,B;x)\bigr)dx\\
 &+\int_1^\infty x^{C-1}
                    \operatorname{Li}_A(e^{-x})\operatorname{Li}_B(e^{-x})dx,\\
 X(A,B,C)={}&\frac{\Gamma(1-A)\Gamma(1-B)}{A+B+C-2}
  +\frac{\Gamma(1-A)\zeta(B)}{A+C-1}
  +\frac{\Gamma(1-B)\zeta(A)}{B+C-1}.
\end{align*}
Jonqui\`ere's expansion at zero \cite[Eq.~25.12.12]{DLMFPolylog} shows that the subtracted integrand
and each fixed parameter derivative are dominated by an integrable
power times a polynomial in $|\log x|$ on a sufficiently small common
polydisc. The tail is exponentially decreasing. Thus $I$ is holomorphic.
Integrating the six subtractions explicitly in the classical Mellin
representation, then using \eqref{qt:polar}, gives
\begin{align}\label{qt:R-integral}
 R(A,B,C)={}&G(C)\bigl(I(A,B,C)+X(A,B,C)\bigr)
            +\frac{G(C)-1}{C}\zeta(A)\zeta(B)\notag\\
 &-\frac{\Gamma(1-A)G(C)-1}{A+C}\zeta(B-1)
  -\frac{\Gamma(1-B)G(C)-1}{B+C}\zeta(A-1).
\end{align}
Every displayed quotient is removable at the origin. For example,
$\Gamma(1-A)G(-A)=1$ proves divisibility by $A+C$.

The derivative $I_{AAB}(0)$ is exactly $J_{21}$, since differentiating
the six subtractions gives $S_{21}$. Likewise
$I_{ABC}(0)+gI_{AB}(0)=K_{11}$, since $G'(0)=g$ and the mixed
subtraction is $S_{11}$. The remaining derivatives in
\eqref{qt:R-integral} are finite Gamma and zeta algebra. In particular,
\[
 \left.\partial_A\partial_C
       \frac{\Gamma(1-A)G(C)-1}{A+C}\right|_0
       =\frac{g^3-Z_3}{3}.
\]
For $C=0$, the derivatives of
$(\Gamma(1-A)-1)/A$ of orders $1,2$ are $\mu_2/2,\mu_3/3$.
These evaluations and the elementary denominators in $X$ give exactly
$E_{21}$ and $E_{111}$.

Finally $f_2f_1-S_{21}=O(x(1+|\log x|^2))$ and
$f_1^2-S_{11}=O(x(1+|\log x|))$ near zero. Their displayed
integrals therefore converge absolutely in the ordinary sense.
\end{proof}

\subsection{Exact cancellation families}

\begin{corollary}[Cyclic quartic reduction]\label{qt:cyclic}
Suppose all pairwise sums of $a,b,c$ are nonzero, and let
$s_1=a+b+c$, $s_2=ab+bc+ca$, $s_3=abc$. Then
\begin{align}\label{qt:cyclic-eq}
 \sum_{\mathrm{cyc}} W_{a,b,c}^{(4)}(0)={}&s_1s_3\Omega_4
 +(-4s_1^2s_2+8s_2^2+12s_1s_3)Ld_3\notag\\
 &+(12s_2^2-36s_1s_3)d_2^2\notag\\
 &+(-2s_1^4+4s_1^2s_2-2s_2^2+24s_1s_3)d_4.
\end{align}
Both off-diagonal coordinates and $e_4$ cancel identically.
\end{corollary}
\begin{proof}
The $\alpha$ coefficient is
$4\sum_{\rm cyc}c(a^3+b^3)-4\sum_{\rm cyc}c^3(a+b)=0$.
The $\beta$ coefficient is $12abc\sum_{\rm cyc}(a+b-2c)=0$.
For the $e_4$ coefficient, the terms with denominator $a+c$
combine to $-b^4$ and cancel the sum of fourth powers. The remaining
symmetric polynomials yield \eqref{qt:cyclic-eq}.
\end{proof}

There is also a useful one-parameter source of completely evaluated
differences. On the additive plane $c=a+b$, the nonclassical part of
\eqref{qt:all-rays} is
\[
 ab(a+b)^2(\Omega_4-12\alpha-12\beta).
\]
Any two admissible rays on this plane therefore have a weighted
difference involving only ordinary zeta derivatives. A small example is
\begin{equation}\label{qt:short-difference}
 \boxed{W_{1,2,3}^{(4)}(0)-\frac92 W_{1,1,2}^{(4)}(0)
 =-20Ld_3+24d_2^2-76d_4+\frac{12}{5}e_4.}
\end{equation}
This is an evaluated identity with no undetermined Tornheim coordinate.

For a single ray with $abc\ne0$, the two displayed off-diagonal
coefficients vanish simultaneously if and only if $a=b=c$.
Indeed the $\beta$ coefficient forces $a+b=2c$; substituting this
into the $\alpha$ coefficient gives a nonzero multiple of
$(a-b)^2$. Thus the nontrivial cubic cancellation conic from the
incoming report does not persist as a conic at the fourth layer.
This assertion concerns the stated formal-coordinate elimination.
An additional arithmetic relation among the constants could enlarge
the family of values reducible in a chosen constant algebra.

\subsection{All-order jet structure and the limit of cyclic diagonal reduction}

The quartic formula is the last layer at which the axis and Euler
stuffle constraints force a cyclic sum to depend on only one diagonal
coordinate. The following exact polynomial statement explains why.

Fix an integer $n\ge3$. Subtract two possible degree-$n$ homogeneous
jets of \eqref{qt:polar} whose axis and Euler boundary restrictions
agree. Their difference is a homogeneous polynomial $U(A,B,C)$
of degree $n$ satisfying
\begin{equation}\label{qt:residual-conditions}
 U(A,B,C)=U(B,A,C),\qquad C\mid U,\qquad
 U(A,0,C)+U(C,0,A)=0.
\end{equation}
These are conditions on a formal reconstruction space, not additional
functional equations asserted for arbitrary perturbations of $T$.

\begin{theorem}[Normal form of the residual jet]\label{qt:all-order-space}
Every polynomial satisfying \eqref{qt:residual-conditions} has a unique
representation
\begin{equation}\label{qt:residual-form}
 U=C\bigl[A(A-C)V(A,C)+B(B-C)V(B,C)\bigr]+ABC\,Q(A,B,C),
\end{equation}
where $V$ is a symmetric two-variable homogeneous polynomial of degree
$n-3$, and $Q$ is homogeneous of degree $n-3$ and symmetric in $A,B$.
Consequently this residual space has dimension
\begin{equation}\label{qt:jet-dimension}
 D_n=\left\lfloor\frac{n^2-1}{4}\right\rfloor.
\end{equation}
For $n=1,2$ the corresponding residual space is zero.
\end{theorem}
\begin{proof}
The boundary polynomial $F(A,C)=U(A,0,C)$ is antisymmetric in $A,C$.
It is divisible by $C$, hence by $A$ as well, and antisymmetry also
implies divisibility by $A-C$. Therefore
$F(A,C)=AC(A-C)V(A,C)$, where $V$ is symmetric and has degree
$n-3$. This uniquely determines $V$. Subtract the first term of
\eqref{qt:residual-form} from $U$. The difference vanishes on both
$A=0$ and $B=0$, by symmetry, and remains divisible by $C$.
Hence it is uniquely $ABCQ$, with the required symmetry.
Conversely the displayed normal form satisfies all three conditions.

The symmetric two-variable polynomials of degree $n-3$ have dimension
$\lfloor(n-1)/2\rfloor$. Polynomials in three variables of that degree
which are symmetric in the first two have dimension
$\lfloor(n-1)^2/4\rfloor$, by counting the exponent pairs of $A,B$.
Their sum is \eqref{qt:jet-dimension}. If $n<3$, the two divisibility
arguments force $U=0$ directly.
\end{proof}

\begin{corollary}[Cyclic residual dimension]\label{qt:cyclic-dimension}
The cyclic symmetrization of the residual space is precisely
\[
 \left\{ABC\,S(A,B,C): S\text{ is symmetric in all three variables
 and homogeneous of degree }n-3\right\}.
\]
Its dimension is
\begin{equation}\label{qt:cyclic-count}
 p_{\le3}(n-3)=\left\lfloor\frac{n^2+3}{12}\right\rfloor,
\end{equation}
where $p_{\le3}(m)$ is the number of partitions of $m$ into at most
three parts. In particular this dimension is one at orders three and
four, and is two at order five.
\end{corollary}
\begin{proof}
The terms involving $V$ in \eqref{qt:residual-form} cancel in pairs
under cyclic summation, because $V(A,C)=V(C,A)$. The remaining sum is
$ABC\sum_{\rm cyc}Q(A,B,C)$. Since $Q$ is symmetric in its first
two variables, this sum is fully symmetric. Conversely any symmetric
$S$ occurs by taking $Q=S/3$ and $V=0$. Symmetric homogeneous
polynomials of degree $m$ have basis $s_1^i s_2^j s_3^k$ with
$i+2j+3k=m$. Counting these triples, or expanding
$(1-t)^{-1}(1-t^2)^{-1}(1-t^3)^{-1}$, gives
$p_{\le3}(m)=\lfloor(m^2+6m+12)/12\rfloor$, proving
\eqref{qt:cyclic-count}.
\end{proof}

Evaluation on the diagonal is a nonzero linear functional on this
cyclic residual space. It therefore determines the whole space when
the dimension is one, explaining the complete cyclic formulas at
orders three and four. The following explicit fifth-order polynomial
shows what fails next:
\begin{equation}\label{qt:fifth-obstruction}
 U_5(A,B,C)=ABC\bigl[(A+B+C)^2-3(AB+BC+CA)\bigr].
\end{equation}
It satisfies the boundary conditions, vanishes on the entire diagonal,
and has nonzero cyclic sum $3U_5$. Consequently the axis, symmetry,
Euler stuffle, and diagonal data alone cannot force a fifth-order
cyclic formula involving just the diagonal fifth derivative and
ordinary zeta jets. This is a precise limitation of those input
identities; other functional equations may provide additional
constraints.

\subsection{Verification and remaining arithmetic questions}

The companion implementation \texttt{check\_quartic.py} checks
the full rational-function identity, cyclic elimination, diagonal and
Euler-diagonal cases, and both elementary integral correction terms.
It evaluates $\alpha$ with a finite Stirling expansion and computes
$J_{21},K_{11}$ by an integrated Jonqui\`ere series on $(0,1)$ plus
exponentially convergent series on $(1,\infty)$. A separate evaluator
computes the fourth ray derivative by a fifteen-node symmetric
difference of the original Mellin continuation. These routes give
\begin{align*}
 \alpha&=13.1294216474884019236329146797203\ldots,\\
 \beta&=6.90381599352243710775587231649323\ldots,\\
 \varepsilon&=4.68913245627221585695112599939183\ldots,\\
 \Omega_4&=798.777376093189588968747119645501\ldots.
\end{align*}
At 55 decimal digits of working precision, four independently evaluated
rays $(1,1,1)$, $(1,0,2)$, $(1,2,2)$, and $(1,-2,3)$ had absolute discrepancies
below $8.2\cdot10^{-34}$. These decimal checks use mpmath and are not
interval certificates; the exact identities rest on the proofs above.

The result completes the quartic finite-coordinate problem posed in the
incoming reports. It does not provide a finite expression for
$\Omega_4,\alpha,\beta$ in a smaller algebra of ordinary constants.
Natural next questions are: relate the two new integral coordinates to
higher harmonic Stieltjes coefficients; identify the dimension and
canonical anchor directions at every order; obtain general cyclic
elimination laws; and formalize the finite jet algebra separately from
the Mellin continuation and normal-convergence lemmas.



% Begin sections/06_audit_research
\section{Audit, interpretation, and further research}
\label{sec:audit}

\subsection{What has been proved, and what has not}

The following table records the mathematical status at the inspected
revision. The word ``resolved'' refers to the explicitly stated component
of a research question, not to every possible arithmetic refinement.

\begin{longtable}{>{\raggedright\arraybackslash}p{0.25\textwidth}p{0.65\textwidth}}
\toprule
Question or object & Result and scope\\
\midrule
\endhead
Nonpolynomial reflected generator &
Theorems~\ref{beta:transform} and \ref{beta:completion} give a
complex-power sine family, its polylogarithmic Euler representation, and
the full joint completion. All Stieltjes indices, all argument-derivative
orders, and all logarithmic powers are covered.\\[5pt]
First nonlinear log-sine bridge &
Theorem~\ref{beta:secondmoment} expresses the sine-square moment through
the inclusive harmonic Stieltjes coefficient and then eliminates that
coefficient in \eqref{beta:cancellation}. The separate cubic bridge is
inherited from \cite{RepoCubic}.\\[5pt]
Centered harmonic cancellation &
Theorem~\ref{hc:main} is necessary and sufficient for every polynomial
in finitely many generalized harmonic variables. It is stronger than
checking finitely many residues. Theorem~\ref{hc:trigamma} evaluates an
infinite family of its regular values.\\[5pt]
Geometric collisions &
Theorems~\ref{gc:main} and \ref{gc:derivative} cover arbitrary finite
clusters and argument-derivative orders at Stieltjes index zero and unit
frequency, with no comparability requirement on the gaps. Higher
Stieltjes indices and iterated merger associativity remain open here.\\[5pt]
Quartic directional Tornheim jets &
Theorem~\ref{qt:main} reconstructs every admissible fourth ray derivative
from three specified coordinates. Proposition~\ref{qt:integrals} gives
convergent representatives. The difference \eqref{qt:short-difference}
eliminates all three coordinates.\\[5pt]
Higher formal relation spaces &
Theorem~\ref{qt:all-order-space} and
Corollary~\ref{qt:cyclic-dimension} give exact dimensions for the stated
axis, symmetry, and Euler boundary system. They do not prove period
independence or completeness of all functional identities.\\[5pt]
Gaussian $S_6$ and revised $S_8$ &
Still conjectural. The present work supplies neither a period certificate
nor a counterexample. Existing formal nonmembership results restrict
their specified search systems only \cite{RepoExact}.\\
\bottomrule
\end{longtable}

The phrase ``finite-coordinate formula'' is deliberately separate from
``evaluation in ordinary constants.'' For example, the quartic theorem
is effective: its retained coordinates are given by convergent Gamma
series and polylogarithmic integrals. A further reduction of those
coordinates is an additional arithmetic theorem. Conversely,
\eqref{beta:cancellation} and \eqref{qt:short-difference} already are
evaluations in the displayed ordinary constants, because the residual
coordinates cancel algebraically.

\subsection{Corrections and qualifications to preserve at integration}

No additional false theorem was established in the focused portions of
the nine incoming reports that were checked. Several previously
documented corrections should retain their original attribution: the
rejected earlier $S_8$ coefficient vector, the restricted sector required
in the Crandall--Jonqui\`ere expansion, the extra harmonic factor in an
external boundary formula, and the distinction between a numerical
period conjecture and formal row-space nonmembership
\cite{RepoExact,RepoCubic,RepoIndependent}. They are not discoveries of
the present article.

The new results do rule out three plausible but unjustified extensions.
They suggest the following concrete wording for future integration.
\begin{enumerate}
\item \textbf{Specify the entire crossing numerator.}
Replace any convention saying only ``take the removable value of the
regularized integral'' by the holomorphic completion
\eqref{beta:Hp}, or its locally equivalent full subtraction. At $p=2$,
the uncompleted analytic limit is $2\zeta(2)$ while the coordinate
finite part is zero. Merely observing that the residue vanishes at the
intersection does not remove its finite contribution.
\item \textbf{Retain boundary contacts under argument differentiation.}
The two derivative correlations satisfy
$C_{(1,0)}+C_{(0,1)}=\pi^2\csc^2(\pi a)$, not zero, in the specified
finite-part convention. A formal integration-by-parts or translation
argument that drops the seam terms loses this exact contribution.
Use the direct repeated-pole calculation or an explicitly compatible
distributional extension.
\item \textbf{Limit the diagonal cyclic claim to its proven orders.}
The polynomial $U_5$ in \eqref{qt:fifth-obstruction} is invisible to the
axis, boundary, and diagonal data but has a nonzero cyclic sum. Thus the
cubic and quartic reduction to one diagonal coordinate cannot be
extended to all orders using those relations alone. Additional
relations must be stated and proved.
\end{enumerate}

There is a related algebraic warning in Section~\ref{hc:section}.
Vanishing at a finite collection of spectral centers is not evidence for
the all-center property. The explicit near-miss polynomials following
\eqref{hc:first-tests} show exactly where this inference fails. The
replacement is the finite polynomial identity $\iota(P)=P$, whose
necessity uses the proved formal independence lemma.

\subsection{Reproducibility and evidence}

Each theorem has an analytic or exact algebraic proof in the article.
The Python files independently check sensitive signs, finite constants,
and symbolic eliminations. They require only Python, SymPy, and mpmath;
no source archive or network connection is needed to run them.
The verification README explains the methods, recorded versions,
precision, and report-writing behavior.

\begin{longtable}{>{\raggedright\arraybackslash}p{0.32\textwidth}p{0.58\textwidth}}
\toprule
Artifact & Independent check\\
\midrule
\endhead
\src{check_beta.py} &
Direct Hurwitz-weight quadrature versus finite Fourier sums and the
Euler--polylog integral; Gamma reflection moments; separately subtracted
digamma logarithmic and cubic integrals; exact Fourier jets.\\[4pt]
\src{centered_checks.py} &
Exact Bernoulli recurrences, affine derivative relations, parity examples,
first-pole counterexamples, and Faulhaber coefficients; independent
finite-head and Hurwitz-tail values for $M=0,\ldots,6$.\\[4pt]
\src{check_quartic.py} &
Exact rational reconstruction, cyclic cancellation, integral correction
constants, dimensions, and the fifth-order obstruction; four fourth
derivatives evaluated from the original continued Mellin expression.\\[4pt]
\src{verify_collision.py} &
Exact quartic and hierarchical coefficient extraction; direct
interval-by-interval finite parts compared with the predicted merged
limit after subtraction.\\[4pt]
\src{verify_derivative_pair.py} &
Independent double-pole quadrature for the two derivative counterterms
and the exact trigamma-reflection sum.\\
\bottomrule
\end{longtable}

The beta identities, except for the integrable singular Euler test,
have observed numerical discrepancies below $1.3\cdot10^{-60}$ in the
recorded run; the Euler test has discrepancy $2.5\cdot10^{-33}$ or less.
The four Tornheim ray comparisons are below $8.2\cdot10^{-34}$.
The derivative-pair reflection checks are below $6.4\cdot10^{-51}$.
These are floating-point observations at the documented precisions.
For geometric collisions, $C-T-Q$ at a finite separation is a remainder
that should tend to zero, not an equality error expected to be tiny at
each separation. The centered checks have an additional explicit
analytic truncation bound proved in the appendix; rounding errors have
not been enclosed by interval arithmetic.

\subsection{Proposed research questions}

The following questions are intended as concrete next projects. They
emphasize exact identities and structure; the analytic estimates needed
to justify an interchange or a numerical reproduction are auxiliary.

\begin{question}[Higher log-sine cancellations]\label{rq:logs}
Determine the smallest harmonic-zeta Laurent family needed for
$J_k(s)$ at each $k\ge3$, using the explicit coefficient formula
\eqref{beta:jet}. Which linear combinations of finite-part products of
digamma, polygamma, and logarithmic sine powers cancel every nonordinary
harmonic coordinate? The first target is the next layer after
\eqref{beta:cancellation}, with its relation matrix derived symbolically
before any numerical search.
\end{question}
The generator fixes both inclusive harmonic conventions and all
endpoint constants. This makes it possible to compare candidate
combinations without changing regularization halfway through the
calculation.

\begin{question}[An arithmetic description of quartic coordinates]
Express $J_{21}$ and $K_{11}$, or their canonical linear combinations,
through Laurent coefficients of explicitly defined harmonic multiple
zeta functions. Does the combination $\beta+\varepsilon$ admit a
shorter one-dimensional harmonic series analogous to the cubic bridge?
Can any of $\Omega_4,\alpha,\beta$ be reduced further in a specified
ordinary constant algebra?
\end{question}
The convergent representations here provide exact starting expressions.
A numerical integer relation can guide a target but would not resolve
the question.

\begin{question}[The second quintic cyclic coordinate]
Compute the full fifth-order cyclic formula. After the known boundary
part is removed, its residual lies in
$ABC\,\mathrm{span}\{(A+B+C)^2,AB+BC+CA\}$.
The diagonal value and one suitable nondiagonal cyclic value, for example
at $(1,1,2)$, separate these two coordinates. Find convergent
representatives and short combinations in which both cancel.
\end{question}
The assertion about the two evaluations is finite algebra: the ratios
$s_1^2/s_2$ are $3$ at $(1,1,1)$ and $16/5$ at $(1,1,2)$.
What remains is to compute the ordinary part and the new analytic
coordinate, not to conjecture the dimension.

\begin{question}[Canonical ray sets at arbitrary order]
Use the normal form \eqref{qt:residual-form} to choose small positive
integer anchor rays giving an invertible evaluation matrix of size
$D_n$. For cyclic reconstruction the corresponding size is
$p_{\le3}(n-3)$. Can the inverses and the resulting cancellation
identities be organized uniformly in $n$?
\end{question}
One should also test additional global Tornheim functional equations
against these spaces. They may reduce the dimension of the actual
reconstruction problem; the dimensions in this article are exact only
for their stated input relations.

\begin{question}[Higher Stieltjes geometric collisions]
Replace $\gamma_0^{(r_i)}$ by $\gamma_{m_i}^{(r_i)}$ in
Theorem~\ref{gc:derivative}. The singular local germs now contain
powers of $\log x$. Construct a simultaneous Mellin-parameter
completion whose derivatives give explicit counterterms uniformly for
hierarchical gaps, and prove that the remainder extends through the
total collision.
\end{question}
The existing all-factor spectral completions are relevant input, but
they do not by themselves justify uniform passage through vanishing
geometric gaps. This is the missing analytic statement to isolate.

\begin{question}[Integer-frequency grids and Fourier weights]
Give the complete geometric theorem for
$\prod_i\gamma_0^{(r_i)}(\{p_i x+s_i\})$, with positive integers
$p_i$, optionally multiplied by an analytic periodic weight. At each
merging grid point the local shifts are $s_i/p_i$ and the pole amplitudes
must be scaled in the ambient coordinate. Derive the resulting finite
Fourier-moment contact formulas and their conductor dependence.
\end{question}
The local interpolation proof strongly suggests this extension. A
complete statement must inventory all grid intersections and preserve
their coordinate scale; it is not included in the proved theorems here.

\begin{question}[Compatibility of successive mergers]
Characterize subtraction rules for partial clusters whose iterated
mergers agree with the simultaneous full subtraction. For three and
four factors, calculate the finite difference between two merger orders
and decide whether a universal local adjustment makes the operation
associative.
\end{question}
The hierarchical example \eqref{gc:hierarchy} supplies an exact test
case. Existence of the fully subtracted limit does not settle
associativity of other subtraction prescriptions.

\begin{question}[Centered products beyond a trigamma square]
Evaluate all centered even moments of products of an even number of
even-order harmonic tails. Start with
$(H_n^{(2r)}-\zeta(2r))(H_n^{(2q)}-\zeta(2q))$, then four tails.
Determine which families reduce by finite summation by parts to
ordinary Euler sums and which require higher-depth constants.
\end{question}
The pole cancellation is already completely answered by
Theorem~\ref{hc:main}; the new task is evaluation of the regular values.
This separates a solved analytic classification from an arithmetic
problem that can grow with the number of factors.

\begin{question}[Other outer shifts and periodic harmonic data]
Classify the polynomials $P$ for which
$\sum_{n\ge0}P(H_n^{(1)},\ldots,H_n^{(R)})(n+a)^{-s}$ is regular
at all nonpositive even integers when $a\ne1/2$.
Then consider harmonic variables with rational inner shifts or periodic
coefficients. What replaces the reflection involution, and how do
coincident translation-pole orbits affect the independence argument?
\end{question}
The additive-difference lemma reduces this to a concrete analysis of
rational telescoping in each pole orbit. It should be proved for the
specific family rather than inferred from a few vanished Bernoulli
coefficients.

\begin{question}[A joint beta--collision generator]
Combine the complex-power sine weight with moving periodic Stieltjes
factors. Determine the polar arrangement in the weight, spectral, and
gap parameters, and give a canonical subtraction whose restrictions
recover both \eqref{beta:Hp} and \eqref{gc:Tr}.
\end{question}
Such a formula would connect Fourier coefficients, Gamma primitives,
and geometric contact terms in a single analytic object. It would also
identify explicitly which order of restriction is compatible with the
unit-coordinate finite part.

\begin{question}[Return to the Gaussian period conjectures]
Find a functional identity or period certificate for the frozen $S_6$
and revised $S_8$ candidates outside the already obstructed relation
system. Can a harmonic or Stieltjes transform supply a new exact
relation after a justified change of variables to the Gaussian alphabet?
\end{question}
No such bridge is proved here. The source vectors and their conventions
should remain fixed, and a proposed proof must produce an exact
certificate for those vectors rather than a nearby fitted identity.

\begin{question}[Formalization in ProveIt]
Separate the formal development into reusable analytic and algebraic
lemmas: finite asymptotic subtraction with normal convergence;
coefficient extraction from a completed meromorphic germ; confluent
interpolation on a fixed contour; the additive-difference lemma; and
the homogeneous residual-space normal form. Which of the present
finite identities can then be generated and checked automatically from
these lemmas?
\end{question}
The supplied exact computations can serve as regression examples. They
are not proof-assistant certificates, and incorporation into the
repository should not silently confer a formalized status.

\appendix
\section{Analytic control of the centered verification tail}
% Begin sections/07_verification_bound
\label{hc:verification-bound}

The numerical checks of the centered trigamma-square moments use a finite
head and an explicitly summed Bernoulli tail. The tail truncation admits
the following analytic bound. This bound controls truncation of an
asymptotic expansion; the recorded \texttt{mpmath} evaluations use ordinary
floating-point arithmetic and are not claimed to be interval enclosures.
The mathematical identity is proved separately in
Theorem~\ref{hc:trigamma}.

For real $x>0$, write
\[
 T(x)=\psi'(x+1/2),\qquad
 A_K(x)=\sum_{j=0}^{K}B_{2j}(1/2)x^{-2j-1},\qquad
 C_K=\frac{2(2K+2)!}{(2\pi)^{2K+2}}.
\]
For every integer $K\ge0$,
\begin{equation}\label{hc:T-remainder}
 0<T(x)\le x^{-1},\qquad
 |T(x)-A_K(x)|\le C_Kx^{-2K-3}.
\end{equation}
Indeed the standard trigamma integral gives
\[
 T(x)=\int_0^\infty e^{-xt}g(t)\,dt,
 \qquad g(t)=\frac{t}{2\sinh(t/2)}.
\]
The partial fraction expansion of the hyperbolic cosecant
\cite[Eq.~4.36.5]{DLMFHyperbolic} is
\[
 g(t)=1+2\sum_{k=1}^{\infty}
             (-1)^k\frac{t^2}{t^2+(2\pi k)^2}.
\]
Expand each rational term through degree $2K$. With $a_k=2\pi k$,
the remainder is exactly
\[
 2(-1)^K t^{2K+2}\sum_{k=1}^{\infty}
              \frac{(-1)^k}{a_k^{2K}(t^2+a_k^2)}.
\]
The positive summand magnitudes decrease in $k$, so the alternating
series estimate yields
\[
 \left|g(t)-\sum_{j=0}^{K}
       \frac{B_{2j}(1/2)}{(2j)!}t^{2j}\right|
 \le \frac{2t^{2K+2}}{(2\pi)^{2K+2}}.
\]
The finite polynomial is identified by the Bernoulli generating
function $g(t)=te^{t/2}/(e^t-1)$. Integration proves the second
inequality in \eqref{hc:T-remainder}; the first follows from
$0<g(t)\le1$.

Fix integers $M\ge0$, $K\ge M$, and $N\ge1$, and set $a=N+1/2$.
Let
\[
 b_j=B_{2j}(1/2),\qquad
 d_\ell^{(K)}=
 \sum_{\substack{0\le j\le K\\0\le\ell-j\le K}}
                    b_jb_{\ell-j},
 \qquad
 C_A=\sum_{j=0}^{K}|b_j|a^{-2j}.
\]
The nondecaying polynomial to subtract from $x^{2M}T(x)^2$ is
\[
 P_M(x)=\sum_{\ell=0}^{M-1}
                  d_\ell^{(K)}x^{2M-2\ell-2}.
\]
This polynomial is independent of $K$ in the stated range; the sum is
empty when $M=0$. Define the finite computational expression
\begin{align}\label{hc:finite-check}
 Q_{M,N,K}={}&
 \sum_{n=0}^{N-1}
   \bigl\{(n+1/2)^{2M}T(n+1/2)^2-P_M(n+1/2)\bigr\}\notag\\
 &+\sum_{\ell=M}^{2K}
          d_\ell^{(K)}\zeta(2\ell+2-2M,a).
\end{align}
Every Hurwitz zeta argument in the last sum is at least $2$.
The ordinarily convergent moment $V_M$ satisfies
\begin{equation}\label{hc:finite-check-bound}
 |V_M-Q_{M,N,K}|
 \le C_K(1+C_A)\zeta(2K+4-2M,a).
\end{equation}
To verify this, for $x\ge a$ use $|A_K(x)|\le C_A/x$ and
\eqref{hc:T-remainder} to obtain
\[
 x^{2M}|T(x)^2-A_K(x)^2|
 \le C_K(1+C_A)x^{2M-2K-4}.
\]
Summing this inequality over $x=n+1/2$, $n\ge N$, proves
\eqref{hc:finite-check-bound}. Expanding the finite square $A_K^2$
and subtracting $P_M$ gives exactly the Hurwitz zeta tail in
\eqref{hc:finite-check}.

The accompanying script uses $N=64$, $K=30$, and $105$ decimal
working digits. Its recorded discrepancies from the exact Bernoulli
formula and the analytic truncation envelopes are as follows:
\begin{center}
\begin{tabular}{rll}
\hline
$M$ & Observed absolute discrepancy & Analytic truncation envelope\\
\hline
0 & $9.911754463\,10^{-80}$ & $1.015061650\,10^{-79}$\\
1 & $4.203540245\,10^{-76}$ & $4.304746532\,10^{-76}$\\
2 & $1.784489779\,10^{-72}$ & $1.827411085\,10^{-72}$\\
3 & $7.583667641\,10^{-69}$ & $7.765877807\,10^{-69}$\\
4 & $3.226610534\,10^{-65}$ & $3.304045785\,10^{-65}$\\
5 & $1.374533946\,10^{-61}$ & $1.407480267\,10^{-61}$\\
6 & $5.863410357\,10^{-58}$ & $6.003761793\,10^{-58}$\\
\hline
\end{tabular}
\end{center}
The script also verifies exactly the first centered Bernoulli pole
test and its stated near-miss example, the formal translation recurrence
through degree $14$, the affine derivative relations for orders $2$
through $8$, nine Faulhaber block constants, and the exact moment
coefficients through $M=10$. These finite checks audit the formulas;
they do not substitute for the all-order proofs.

% Primary source for the partial fraction identity:
% https://dlmf.nist.gov/4.36#E5



% Begin references
\begin{thebibliography}{99}
\small
\setlength{\itemsep}{3pt}
\setlength{\parskip}{0pt}
\addcontentsline{toc}{section}{References}

\bibitem{RepoExact}
ProveIt incoming report,
\emph{Exact Resonant Identities Beyond the Diagonal}, 11 October 2026.
\href{https://github.com/VladimirReshetnikov/ProveIt/blob/c2cb285b64273e4cb2c8d7a03c68a4a99ac812d4/docs/incoming/ProveIt_Exact_Resonant_Identities_2026-10-11.zip}{Pinned source archive}.
Especially Sections 2--4 and 6--7, Questions 8 and 11.

\bibitem{RepoIndependent}
ProveIt incoming report,
\emph{Independent Orders and Exact Reductions in Polylogarithm--Zeta
Calculus}, 11 October 2026.
\href{https://github.com/VladimirReshetnikov/ProveIt/blob/c2cb285b64273e4cb2c8d7a03c68a4a99ac812d4/docs/incoming/ProveIt_Independent_Orders_2026-10-11.zip}{Pinned source archive}.
Independent Tornheim orders, boundary jets, and Section 7.

\bibitem{RepoCubic}
ProveIt incoming report,
\emph{Cubic Tornheim Derivatives and Harmonic Stieltjes Constants},
11 October 2026.
\href{https://github.com/VladimirReshetnikov/ProveIt/blob/c2cb285b64273e4cb2c8d7a03c68a4a99ac812d4/docs/incoming/ProveIt_Cubic_Harmonic_Mixed_Orders_2026-10-11.zip}{Pinned source archive}.
Especially the cubic harmonic bridge and Section 7, Question 4.

\bibitem{RepoPeriodic}
ProveIt incoming report,
\emph{Periodic Stieltjes Collisions: Exact Contact Terms, Finite-Part
Fubini Laws, and Gamma--Polylogarithm Primitives}, 10 October 2026.
\href{https://github.com/VladimirReshetnikov/ProveIt/blob/c2cb285b64273e4cb2c8d7a03c68a4a99ac812d4/docs/incoming/ProveIt_Periodic_Collision_Contacts_2026-10-10.zip}{Pinned source archive}.

\bibitem{RepoContact}
ProveIt incoming report,
\emph{Periodic Stieltjes Contact Calculus: All-Index Collision
Corrections, Harmonic Derivative Anomalies, and Exact Zeta-Jet
Integrals}, 10 October 2026.
\href{https://github.com/VladimirReshetnikov/ProveIt/blob/c2cb285b64273e4cb2c8d7a03c68a4a99ac812d4/docs/incoming/ProveIt_Periodic_Stieltjes_Contact_2026-10-10.zip}{Pinned source archive}.

\bibitem{RepoMixed}
ProveIt incoming report,
\emph{Mixed Spectral Identities: Harmonic Newton Series, Dougall Jets,
Unequal Twists, and Stieltjes Collisions}, 10 October 2026.
\href{https://github.com/VladimirReshetnikov/ProveIt/blob/c2cb285b64273e4cb2c8d7a03c68a4a99ac812d4/docs/incoming/ProveIt_Mixed_Spectral_Identities_2026-10-10.zip}{Pinned source archive}.
All-factor spectral completion and the three-factor contact
specialization; internal paths are recorded in \src{SOURCE_AUDIT.md}.

\bibitem{EspinosaMoll}
O.~R. Espinosa and V.~H. Moll,
\emph{On some definite integrals involving the Hurwitz zeta function},
Ramanujan Journal \textbf{6} (2002), 159--188.
\url{https://arxiv.org/abs/math/0012078}.
The cited section numbers refer to the 34-page arXiv version.

\bibitem{Kargin}
L. Karg\i n, A. Dil, M. Cenkci, and M. Can,
\emph{On the Stieltjes constants with respect to harmonic zeta functions},
2023. \url{https://arxiv.org/abs/2304.03517}.
Related journal DOI: \url{https://doi.org/10.1016/j.jmaa.2023.127302}.

\bibitem{Can}
M. Can, L. Karg\i n, M. Cenkci, and A. Dil,
\emph{On certain harmonic zeta functions}, 2024.
\url{https://arxiv.org/abs/2403.07123}.
Version 1, especially its introductory harmonic-zeta Laurent
normalizations.

\bibitem{BorweinDilcher}
J.~M. Borwein and K. Dilcher,
\emph{Derivatives and fast evaluation of the Witten zeta function},
Ramanujan Journal \textbf{45} (2018), 413--432.
\url{https://doi.org/10.1007/s11139-017-9890-9}.
\href{https://carmamaths.org/resources/jon/omega3.pdf}{Author manuscript}.

\bibitem{HardouinSinger}
C. Hardouin and M.~F. Singer,
\emph{Differential Galois theory of linear difference equations},
Mathematische Annalen \textbf{342} (2008), 333--377.
\url{https://arxiv.org/abs/0801.1493}.
Context for the additive-difference mechanism; the lemma used here is
proved in full in Section~\ref{hc:section}.

\bibitem{DLMFHurwitz}
NIST Digital Library of Mathematical Functions,
\S25.11, \emph{Hurwitz Zeta Function}.
\url{https://dlmf.nist.gov/25.11}.
Fourier expansion, shift and argument-derivative identities, and
Lerch's formula.

\bibitem{DLMFBeta}
NIST Digital Library of Mathematical Functions,
\S5.12, \emph{Beta Function}.
\url{https://dlmf.nist.gov/5.12}.

\bibitem{DLMFGamma}
NIST Digital Library of Mathematical Functions,
\S5.11, \emph{Asymptotic Expansions}.
\url{https://dlmf.nist.gov/5.11}.

\bibitem{DLMFPolylog}
NIST Digital Library of Mathematical Functions,
\S25.12, \emph{Polylogarithms}.
\url{https://dlmf.nist.gov/25.12}.
In particular Eq.~25.12.12.

\bibitem{DLMFHyperbolic}
NIST Digital Library of Mathematical Functions,
\S4.36, \emph{Infinite Products and Partial Fractions}.
\url{https://dlmf.nist.gov/4.36}.
In particular the hyperbolic-cosecant expansion, Eq.~4.36.5.

\end{thebibliography}


\end{document}
